% 本文件由 LaTeX 排版 Agent 根据有效 translation.json 自动生成。
% author=Ephraim the Syrian; work_id=3703; title=圣诞颂歌

\chapter{圣诞颂歌}

% structure: p0001 source=h1 target=omitted reason=page-title-already-rendered-by-parent

% structure: p0002 source=h2 target=section reason=relative-heading-rank; base=h2

\section{第一首颂歌}

这是令先知、君王和司祭欢欣的日子，因为在这一天，他们的话语都应验了，所有的一切都真正成就了！因为今日童贞玛利亚在\Place{白冷}生下了\Person{厄玛奴耳}。依撒意亚在古时所说的话，\BibleRef{依 10:19}今天成为了现实。祂在那里出生，要记录外邦人的数目！达味曾咏唱的圣咏，今天因着应验而来到！米该亚曾说的话语，\BibleRef{米 5:2}今天确实应验了！因为来自厄弗辣大的一位牧者出现了，祂的牧杖管辖着灵魂。看！由雅各伯闪耀了一颗星，\BibleRef{户 24:17}从以色列兴起了一个首领。\BibleRef{欧 1:11}巴郎所说的预言，今天得到了解释！隐藏的光也降下，从身体中显出祂的美！在匝加利亚身上说话的光，今天在白冷照耀了！\par

王国的光已经升起，在厄弗辣大，君王的城里。雅各伯祝福的祝福，今天应验了！同样，那生命树带给了凡人希望！撒罗满隐藏的箴言\BibleRef{箴 3:18}今天得到了解释！今天诞生了那孩子，祂的名字被称为\Quote{奇妙}！\BibleRef{依 9:6}因为天主以婴孩形态显现，这真是奇妙。圣神以\Quote{虫}一词用比喻预示了祂，因为祂的诞生没有婚姻。圣神所预表的预像，今天其含义得到了【解释】。祂在祂面前如同根芽，如同干地中的根芽。\BibleRef{依 53:2}一切隐密所说的话，今天都公开成就了！那隐藏于犹大的君王，塔玛尔从祂大腿窃取了祂；今天崛起了祂凯旋的美，她曾在隐藏的状态中爱慕祂。卢德躺在波阿次身边，因为她察觉到了隐藏于他内的生命之药。今天她的誓愿应验了，因为从她的后裔中兴起了万物的苏醒者。亚当将产痛带给女人，因为她从他而来。今天她赎回了她的产痛，因为她生了拯救她的救主！我们母亲厄娃生了一个男人，而祂自身没有受生。厄娃的女儿如何更应被相信没有男人就生了孩子！处女地生下了那个亚当，他是大地的元首！今天童贞女生下了那亚当，他是诸天的元首。亚郎的棍杖发了芽，干木结出了果实！它的奥迹今天显明了，因为童贞的子宫生了一个孩子！\par

那相信先知的话语的人群蒙受了羞辱；因为除非我们的救主来了，他们的话语就被证实为虚假！愿那来自真理之父的\Quote{真实者}受赞美，祂实现了真实先见者们的话语，这些话语在他们的真理中应验了。主啊，请从祢的宝库，从祢圣经的库房中，拿出古时义人的名字，他们曾切望看到祢的来临！代替亚伯尔的舍特，预示了圣子被杀害；因祂的死亡，加音带入世界的嫉妒被减弱了！诺厄看见了天主的儿子们——圣者们突然变得放荡，他盼望那神圣的圣子，通过祂，淫荡之人被转变为圣洁。那遮盖诺厄的两个兄弟，\BibleRef{创 9:23}看见了天主的独生子将要来遮盖亚当的裸体，他被骄傲沉醉。闪和耶斐特，因着恩慈，盼望那恩慈的圣子，祂要来使客纳罕从罪的奴役中得释放。\par

默基瑟德期待祂；作为祂的代表，他盼望能看到司祭职的主，祂的牛膝草\BibleRef{肋 14:52}净化世界。罗特看见索多玛人如何歪曲本性：他盼望那本性之主，祂赐予非本性的圣洁。亚郎寻找祂，因为他看见，如果他的棍杖吞掉了蛇，\BibleRef{出 7:12}那么祂的十字架将吞掉那吞吃了亚当和厄娃的古蛇。梅瑟看见了那举起的铜蛇，它治愈了毒蛇的咬伤，他盼望看到那医治古蛇创伤的祂。梅瑟看见他独自一人从天主那里保持了光芒，他盼望那将来并因祂的教导\OriginalTerm{使众神增多}{multiplied gods}的祂。\par

探子加肋布将一串（葡萄）挂在棍杖上，来到并切望看到那串（葡萄），祂的酒将慰藉世界。农的儿子若苏厄渴望祂，以便他能领会他自己姓氏的力量：因为因祂的名，他变得如此强大，\BibleRef{希 4:8}何况祂藉着诞生本身呢？这位若苏厄采集、携带并带来了一些果实，他渴望生命树，要品尝那使万物苏醒的果子。辣哈布也寻找祂；因为当朱红绳子以预像将她从愤怒中救赎时，她以预像尝到了真理。厄里亚渴望祂，当他在地上没有看见祂时，他通过彻底洁净的信德，升天去看见祂。梅瑟和厄里亚看见了祂；那温顺的人从深处上升，那热忱者从高处下降，并在中间看见了圣子。他们预表了祂来临的奥迹：梅瑟是死者的预像，厄里亚是生者的预像，他们在祂来临时飞翔迎接祂。\BibleRef{得前 4:17}因为那些已尝过死亡的死者，祂使他们最先；其余未被埋葬的，最后被提升去迎接祂。\par

有谁能为我数算那些期待圣子的义人？他们的人数是我们软弱的受造物之口所不能确定的。请为我祈祷，亲爱的，好使我在另一时日得着力量，在另一篇颂歌中，尽我所能，如此陈述他们的预尝。谁能堪当颂扬那为我们而升起的真理之子？因为义人们长久渴慕祂，盼望在自己的世代得以看见祂。亚当期待祂，因为祂是革鲁宾的主，能赐予进入和安居在生命树的枝条之旁。亚伯尔渴慕祂，愿祂在他有生之日来临，好使他能看见天主的羔羊，替代他所献的那只羔羊。厄娃也期待祂，因为妇女的赤身露体是难堪的，而祂有能力给他们披上衣服，不是用树叶，而是用他们所失去的那同一光荣。那由多人建造的塔，在奥秘中期待一位，祂要降临，在地上建造一座高耸入天的塔。是的，活物的方舟以预像期待我们的主，因为祂要建造圣教会，在其中灵魂寻得避难之所。在培肋格的日子，大地被分作七十种语言。祂要藉着语言将大地分给祂的宗徒。曾被洪水吞没的大地，在沉默中向她的主呼号。祂降临并开启了圣洗之门，人们被它吸引升天。舍特、厄诺士和刻南都被称为天主的儿子；他们期待天主子，好使他们因恩宠成为祂的弟兄。默突舍拉活了将近一千年：他期待那使继承者得享永生的圣子！恩宠本身在隐藏的奥秘中为他们恳求，愿他们的主在他们的时代来临，补足他们的缺失。因为圣神在他们内，代替他们以默祷恳求：\BibleRef{罗 8:26}祂激动了他们，在他们内他们仰望那位他们所渴慕的救赎者。\BibleRef{伯前 1:11}\par

义人的灵魂在圣子身上看出生命之药；因此它渴望祂能在自己的有生之日来临，好品尝祂的甘饴。哈诺客渴慕祂，既然在地上他未曾看见圣子，他便因伟大的信德而成义，升到天上见祂。有谁会拒绝恩宠呢？当古人历尽艰辛未能获得的恩赐，如今白白地赐予世人时。拉默客也期待祂来临，仁慈地使他从劳苦和双手的操劳中得享安息，并从那义者所诅咒的地上得享安息。\BibleRef{创 5:29}于是拉默客看到自己的儿子诺厄——在他身上，有关圣子的预像被刻画出来。代替遥远的主，近在眼前的预像带来了安息。是的，诺厄也渴望见到祂，他己尝过祂助佑恩宠的滋味。因为如果祂的预像保存了活物，祂本人岂不更能赐予灵魂生命！诺厄渴慕祂，通过考验认识祂，因为藉着祂方舟得以建立。因为如果祂的预像如此拯救了生命，祂本人无疑会救得更多。亚巴郎在圣神中感知圣子的诞生尚遥远；他因见到祂的日子而喜乐，以代替亲见祂本人。\BibleRef{若 8:56}依撒格渴望见祂，因为他尝过了祂救赎的滋味；\BibleRef{希 11:19}因为如果祂的标记如此赋予生命，祂本人藉着实体更会如此。\par

今天守护者们欢乐，\BibleRef{达 4:13}因为警醒者来唤醒我们！有谁愿意在这黑夜中沉睡，而整个世界都在警醒？自从亚当因罪恶将死亡的睡眠带入世界，警醒者降临，为将我们从罪恶的沉睡中唤醒。我们不要像放债者那样警醒，他们想着利息，常常夜间警醒，计算本金和利息。盗贼警醒而谨慎，他将睡眠埋藏在地里。他的警醒全在于此：使睡着的人大大警醒。贪食者同样警醒，他吃得过多而辗转不安；他的警醒对他来说成为折磨，因为他不能忍受节制。商人同样警醒，夜间他活动手指，计算进来的银两，看看财富是否翻倍或三倍。富人同样警醒，他的财富驱走睡眠：他的狗在睡，他却看守财宝以防盗贼。忧虑的人也警醒，他的忧虑吞噬了睡眠：虽然他的末日就在枕边，他却为未来多年的忧虑而醒着。撒殚教导，我的弟兄们，用一种警醒代替另一种警醒：对善行昏昏欲睡，对恶行清醒警醒。连犹达斯依斯加略整夜警醒，出卖了那赎买整个世界的义血。黑暗之子披上黑暗，因为他脱去了光明；那创造白银的，盗贼为了白银出卖了祂。是的，法利塞人，黑暗之子，整夜保持警醒：黑暗之子警醒，为要遮蔽那无限的光明。那么你们要像天上的光体一样，在这星光灿烂的夜晚警醒。因为虽然它的颜色如此黑暗，但在德性上却是明亮的。\par

因为无论谁像这光明者一样，在黑暗中警醒祈祷，在这可见的黑暗中，就有不可见的光围绕他！恶人虽站在日光之下，却像黑暗之子行事；虽然外表披着光明，内里却被黑暗束缚。亲爱的，我们不要因自己正在警醒而受骗！因为凡警醒得不正当的，他的警醒是不义的警醒。凡不喜乐地警醒的，他的警醒不过是沉睡；凡不纯洁地警醒的，他的清醒反成他的仇敌。这是嫉妒者的警醒！一团坚实之物，充满伤害。那种警醒不过是交易，充满轻蔑和嘲弄。若易怒者警醒，他的清醒将充满愤怒和烦躁，他的警醒对他全是怒气与咒骂。若饶舌者警醒，他的口便成为通道，对罪恶敞开，对祈祷却是阻碍。\par

智者若谨守不寐，必择其一：或享甜适之眠，或守神圣之警醒。是夜何其美，因那\Quote{美者}\BibleRef{歌 1:15} 升起来，欲使我们成为美。愿一切扰乱不侵入我们的守夜！愿耳的听觉纯洁，眼的视界贞洁，心的沉思圣化，口的言语清明。玛利亚今日在我们内隐藏了来自亚巴郎的酵母。让我们因此怜悯乞丐，如同亚巴郎怜悯贫困者。今日，从温和的达味家族，凝乳降在我们身上。愿人向迫害者显示仁慈，如同耶西之子对待撒乌耳。\EditorialNote{撒慕尔纪上第二十六章} 先知甘美的盐\BibleRef{列下 2:20}今日倾注在外邦人中。让我们藉那使古人失味之物，获得新味\BibleRef{玛 5:13}。让我们讲说智慧之言，不说其外的事，免得我们自己处于智慧之外！\par

在这和好之夜，愿无人发怒或忧郁！在这安宁之夜，愿无威胁或搅扰！此夜属于甘美者，愿其中无苦涩或严厉！此夜属于温良者，愿其中无高傲或狂妄！在这宽恕之日，愿我们不苛求过失！在这喜乐之日，愿我们不散布悲伤！在这如此甘美之日，愿我们不苛刻！在这安息之日，愿我们不愤怒！在这天主来就罪人之日，愿义人不在心中自高于罪人！在这万有之主来就仆人之日，愿主人也慈爱地俯就仆人！在这富者为我们成为贫穷之日，愿富者与穷人在他的桌旁分享。今日，那恩赐虽未蒙祈求，却已临于我们！因此，让我们向那些向我们哭求的人施舍。此日为我们敞开了高空之门，使我们的祈祷得以进入。让我们也为那些犯过并向我们求恕的恳求者敞开大门。今日，自然之主违反本性而改变；愿我们不以改变恶意为烦扰。身体固于本性，不能变大或变小；但意志拥有如此权能，能增长至任何尺度。今日，天主性封印在人性上，使人性能以天主性的印记为装饰。\par

% structure: p0014 source=h2 target=section reason=relative-heading-rank; base=h2

\section{赞美歌 2}

愿那今日使白冷喜乐的圣婴受赞美！愿那今日使人类复返青春的婴孩受赞美！愿那俯就我们饥饿状态的果实受赞美！愿那突然丰富我们的匮乏、供应我们需要的善者受赞美！愿祂的慈悲使祂屈尊俯就我们软弱的祂受赞美！\par

赞美那为我们赎罪而遣来的泉源。赞美那藉成全安息日而使其无效的祂！赞美那斥责癞病、癞病即消的祂，那热病一见即逃的祂！赞美那担当我们劳苦的仁慈者！光荣归于祢的来临，祢使人类苏醒！\par

光荣归于那位藉首生子来到我们中间的祂！光荣归于那藉其声音发言的寂静！光荣归于那藉其晨光显现的高天者！光荣归于那乐意拥有身体的属神者，为使人能感受其德能，并藉那身体向祂家中的身体施以怜悯！\par

光荣归于那隐藏者，其子已显明！光荣归于那永生者，其子已受死！光荣归于那伟大者，其子降下而成为微小！光荣归于那权能，祂以形状约束了自己的伟大，以形象约束了自己不可见的本性！以眼和心，我们目睹了祂，是的，以两者。\par

光荣归于那隐藏者，即使用心智，探求者也完全不能感知祂；但藉祂的恩慈，祂被人手感知了！那不可触摸的本性，被祂的手捆绑束缚，被祂的脚刺穿举起。祂甘愿自取身体，交予那些捉拿祂的人。\par

愿那被自由意志钉死、因祂容许了的祂受赞美！愿那被木头承载、因祂准许了的祂受赞美！愿那被坟墓束缚、因而为其设限的祂受赞美！愿那自愿进入母胎与诞生、怀抱与成长的祂受赞美！愿那以变化为人性购得生命的祂受赞美！\par

愿那印封我们灵魂、装饰并聘娶了灵魂的祂受赞美！愿那使我们的身体成为祂不可见本性之帐幕的祂受赞美！愿那藉我们的口舌诠释了祂奥秘事的祂受赞美！让我们赞美那以我们的琴弦歌颂其光荣、以竖琴歌颂其德能的圣言。外邦人已聚集而来，聆听祂的诗章。\par

光荣归于善者的圣子，被恶者之子所弃绝！光荣归于义者的圣子，被邪恶之子所钉死！光荣归于那释放我们、却为我们被捆绑的祂！光荣归于那交付保证、又赎回保证的祂！光荣归于那美丽的祂，使我们符合祂的肖像！光荣归于那俊美的祂，不顾我们的丑陋！\par

光荣归于那在黑暗中播撒其光明、在隐藏中受辱、并遮盖其奥秘的祂。祂也脱去并剥夺了我们污秽的衣服。\BibleRef{匝 3:3} 愿光荣归于高天的祂，祂将祂的盐\BibleRef{谷 9:49}调和在我们的心智中，将祂的酵母调和在我们的灵魂中。祂的身体成为饼，以苏醒我们的死寂。\par

赞美那为我们众人偿付了祂未曾借债的富者；祂写下帐单，又成了我们的债务人！藉祂的轭，祂为我们折断那俘掳我们的锁链。光荣归于那受审判的审判者，祂使十二位宗徒坐在审判座上审判各支派，并藉无知者定那国经师的罪！\par

光荣归于那永不能为我们所衡量的祂！我们的心太小，不足以容纳祂，我们的心智太弱。祂以智慧的财富使我们的小见识显得愚拙。光荣归于那俯就并询问的祂\BibleRef{路 2:46}，为能听见并学习祂所知道的事；为能藉问题启示祂有益恩宠的宝藏！\par

让我们朝拜祂，祂以祂的教训光照了我们的心智，在我们的听觉中为祂的言语寻找了途径。赞美祂，祂将果实嫁接在我们的树上。感谢那位派遣祂的嗣子的祂，为藉祂吸引我们归向祂，使我们与祂同为嗣子！感谢那位善者，一切善的起因！\par

有福的那一位，祂未曾责骂，因为祂是良善的！有福的那一位，祂未曾嫌弃，因为祂也是公义的！有福的那一位，祂沉默不语，却又责备，好藉着两者使我们复生！祂的沉默是严厉的，且带有斥责。祂的严厉是温和的，即便在控告时也是如此；因为祂责备了叛徒，却亲吻了盗贼。\par

光荣归于我们理智的隐藏农夫！祂的种子落入了我们的田地，使我们的心灵富足。祂的出产百倍地进入了我们灵魂的宝库。让我们朝拜祂，祂坐下休息，又在路上行走，好使\Quote{道路}行在路上，也使\Quote{门}为那些进入的人预备，藉以进入天国。\par

有福的牧人，祂为我们的和好而成为羔羊！有福的枝条，祂成了我们救赎之杯！有福的葡萄串，生命之药的泉源！有福的农夫，祂成为麦子，好被播种；成为禾捆，好被收割！\newline{}【有福的】建筑师，祂成为高塔，作我们的安全之所！有福的那一位，祂如此调节我们心灵的感受\BibleRef{箴 18:10}，使我们能用竖琴唱出那带翼的受造物以其曲调所无法歌唱的！光荣归于祂，祂看见我们如何喜欢在愤怒和贪婪中如同畜类，便降临并成为我们中的一员，好使我们能成为属天的！\par

光荣归于祂，祂从不缺乏我们的赞美，却因仁慈待我们而感到需要；祂因爱我们而饥渴\BibleRef{玛 25:40}，并请求我们给祂，又渴望给我们。祂的果实与我们人类相融，好使我们在祂内能亲近那屈尊就卑的祂。藉着祂树干上的果实，祂将我们接枝到祂的树上。\par

让我们赞美祂，祂得胜了，并藉着祂的鞭伤使我们复生！赞美祂，祂藉着祂的荆棘除去了诅咒！赞美祂，祂藉着祂的死亡将死亡置于死地！赞美祂，祂保持沉默，却使我们成义！赞美祂，祂斥责了那曾胜过我们的死亡！有福的祂，祂助佑的恩宠涤净了左边！\par

赞美祂，祂醒着，却使那俘虏我们的陷入沉睡。赞美祂，祂入睡，却驱散了我们的深睡。光荣归于天主，祂医治了软弱的人性！光荣归于祂，祂受了洗，将我们的不义淹溺在深渊中，并扼住了那扼杀我们的\BibleRef{路 8:33}！让我们用所有的口舌赞美万有的主！\par

有福的医生，祂降临，无痛地施行切割，并用并非苛刻的药物医治创伤。祂的圣子成为药物，向罪人显示慈悲。有福的祂，居住在母腹中，在那里建造了一座完美的圣殿，好使祂能居住其中；一座宝座，好使祂能在其中；一件外衣，好使祂能披上它；一件武器，好使祂能藉以得胜。\par

有福的祂，我们的口舌无法充分赞美，因为祂的恩赐远非演说家的技巧所能讲述；诸般官能也无法充分赞美祂的良善。因为我们无论怎样赞美祂，都显得不足。\par

既然沉默和约束自己是无用的，愿我们的软弱能宽恕我们所能献上的这种赞美。\par

祂何其仁慈，从不要求超过我们能力所能给予的！你的仆人若没有献上他所能给的，反而拒绝他所欠的，又怎能逃避本利上的定罪？光荣的海洋，你并不需要你的光荣被歌颂，请以你的良善接纳这一滴赞美；因为藉着你的恩赐，你已赐给我舌头以赞美你的感官。\par

% structure: p0037 source=h2 target=section reason=relative-heading-rank; base=h2

\section{颂歌第三首}

有福的是你的那第一日，主啊，你的庆节之日便铭刻其上！你的日子如同你一样，它向人显示慈悲，它被传递下来，随同所有世代而来。\par

这是一个终结于年长者，又返回以开始于年幼儿的日子！一个因爱而自我更新，好能以其大能更新我们这些朽坏之物的日子。你的日子，当它探访我们并过去之后，又慈悲地回来再次探访我们；因为它知道人性需要它，在一切事上如同你一样，寻求我们。\par

世界需要它的泉源；主啊，对于它，如同对于你一样，其中的一切都感到干渴。这是一个支配季节的日子！你的日子的权柄如同你的权柄一样，伸展到已来的和将来的世代！你的日子如同你一样，因为它虽为一，却萌芽并繁衍自身，好能如同你！\par

主啊，在这你的日子中，它临近我们，我们看见你遥远的诞生！愿你的日子对我们如同你一样，主啊，愿它成为和好的中保和担保。\par

你的日子使天地和好，因为在这一天，至高者降到了至低者。\par

你的日子能使那因我们的罪而发怒的义者与人和好；你的日子赦免了成千的罪过，因为在这一天，怜悯的心肠向有罪者照耀了出来！\par

主啊，你的日子是伟大的，求你不要让它在我们身上变得微小；愿它照它素常所行的，向我们这些悖逆者显示慈悲！\par

主啊，若每一天祢的宽恕都涌流出来，在这日子上，它该是何等伟大！所有的日子都从祢光明日子的宝库中获取祝福。所有庆节都从这庆节的储藏中获得了它们的光彩与装饰。主啊，求祢在这日子上，将祢仁慈的腑脏丰沛地赐予我们！求祢使我们能辨识祢的日子与所有日子不同！因为祢诞生日子的宝库是伟大的；愿它成为债务人的赎价！这日子超越所有日子，是伟大的，因为在这一天，仁慈临到了罪人。祢这伟大的日子是药物的仓库，因为在这一天，生命之药照耀了受伤者！这日子是助人恩宠的宝藏，因为在这一天，光明向我们的盲目闪耀！是的，它也为我们带来了一捆禾穗；它来了，为的是从中流出丰盛以消除我们的饥饿。这日子正是那先驱之串，其中隐藏着救恩之杯！这日子是首生之庆节，它首先诞生，胜过了所有庆节。在冬天，它使不结果子的葡萄树枝上的果实脱落，却有果实为我们生出\BibleRef{依 5:2}；在寒冷中，它使所有树木都光秃，但有一枝嫩芽从\Person{叶瑟}家族为我们发绿。在十二月，当种子隐藏在地里时，生命之穗从母胎中抽芽。在三月，当种子在空中发芽时，一捆禾穗\BibleRef{肋 23:10}将自己播种在地里。那收成，死亡在地狱中吞噬了它；然而隐藏其中的生命之药却将它冲破了！在三月，当羔羊在荒野中咩叫时，逾越节羔羊进入了母胎！从渔夫上来的河流中，祂受洗并上来，祂用祂的网捕捞万物。从\Person{西满}捕到鱼的河流中，人灵的渔夫上来，并捕获了他。祂用捕获一切强盗的十字架，将那个强盗带到生命之中！\BibleRef{路 23:43}永生者以祂的死亡清空了地狱，祂解开了它，放出了其中的整个群众！税吏与娼妓，那不洁的陷阱，狡诈捕鸟者的罗网，圣者一并获取了！那个有罪的女人，曾是为人的陷阱，祂却使她成为忏悔女性的明镜！那棵无花果树，它抛弃了果实，拒绝结果，却献上了\Person{匝凯}作为果实；它没有结出自己本性的果子，却结出了一个理性的果子！主将祂的渴望铺展在井上，用祂向那口渴者所求的水捕获了她。祂在井旁捕获了一个灵魂，又与她一同捕获了整个城市：\BibleRef{若 4:42}圣者捕获了十二位渔夫，又与他们一同捕获了整个世界。至于\Person{依斯加略}，他逃脱了祂的网，但绞索却套在了他的脖子上！祂那赋予生命的网捕获活人\BibleRef{玛 13:47}，而逃避它的人，正逃避生命！\par

谁能告诉我，主，在祢内隐藏的诸多援助呢？干渴的口怎能从天主性的泉源中饮用？\newline{}今日应允我们恳求的声音；让我们言语的祈祷在行动中生效。治愈我们，我主；每次我们看到祢的庆节，愿它使我们所听的流言消逝。我们的心神在这些声音中徘徊。圣父的声音，平息其他的声音；世界喧嚣，在祢内愿它获得宁静；因为藉着祢，大海从风暴中平静。\newline{}魔鬼听到亵渎的声音时欢喜：愿守望者照常因我们而喜乐。\BibleRef{玛 18:10} 在祢的羊群中有悲伤的声音；哦，使一切喜乐的祢，\BibleRef{路 15:7} 愿祢的羊群喜乐！\newline{}至于我们的怨言，我主，不要在其中弃绝我们：我们的口发出怨言，因为它是罪恶的。愿祢的日子，主，赐给我们各种喜乐，以和平的花朵，让我们庆祝祢的逾越节。\newline{}在祢升天之日，我们被高举：\BibleRef{若 20:17} 新饼将成为其纪念。\newline{}主，增加我们的和平，使我们可以庆祝天主性的三个庆节。\newline{}伟大的日子，主，愿我们不被轻视。众人尊崇祢诞生的日子。祢，义者，保持祢诞生的荣耀；因为连黑落德也尊崇他诞生的日子！不洁者的舞蹈取悦了暴君；主，愿贞洁妇女的声音在祢面前甘甜！主，愿贞洁妇女的声音取悦祢，祢圣洁地守护她们的身体。\newline{}黑落德的日子像他自己：祢的日子也像祢！那个烦恼者的日子被罪困扰；祢的日子是多么美丽，正如祢一样！暴君的庆节杀害了宣讲者；在祢的庆节，人人宣讲荣耀。在凶手的日子，声音被缄默；但在祢的日子，是庆节的声音。那污秽者在自己的庆节中熄灭了光明，使黑暗能遮盖通奸者。圣者的季节修剪灯盏，使黑暗与其隐藏之事一同逃遁。那只狐狸的日子\BibleRef{路 13:32} 臭如自己；但真羔羊的庆节是圣洁的。\newline{}恶人的日子像他一样消逝；祢的日子像祢自己永远长存。暴君的日子像他自己一样咆哮，因为用他的锁链它缄默了义人的声音。温良者的庆节像祂自己一样宁静，因为祂的太阳照耀祂的迫害者。\newline{}那暴君意识到祂不是王，因此祂向万王之王让位。主，整日都不足以平衡祢的赞美与他的谴责。愿祢的恩慈之日使我的罪消逝，因为我是用不洁者的日子来称量祢的日子！因为祢的日子远超比较！也不能与我们的日子相比。人的日子如同属土者；天主的日子如同天主！主，祢的日子大于先知们的日子，而我竟然拿来放在凶手的旁边！主，祢知道一切，知道如何听取我的舌头所做的比较。愿祢的日子赐给我们生命的请求，因为他的日子赐予死亡的请求。那贫穷的王在他的庆节上起誓，半壁江山作为舞蹈的奖赏！那么，祢的庆节，哦，使一切富足的祢，愿祢慈悲地洒下一粒细麦面粉！\newline{}从干地涌出泉源，足以满足外邦人的干渴！从童贞女的胎中，如同从坚固的磐石，种子萌芽，结出许多果实！若瑟装满了无数的谷仓；\BibleRef{创 41:49} 但在饥荒的岁月中，它们被倒空并缺乏。一束真正的麦穗赐下食粮；天上的食粮，没有限量。\newline{}长子在旷野中擘开的饼，虽极好，却消逝了。祂再次返回并擘开新饼，这饼时代和世代不会耗尽！祂擘开的七个饼也失败了，\BibleRef{玛 15:36} 祂增多的五个饼也被吃尽；\BibleRef{玛 14:17} 祂擘开的饼超过了世界的需要，因为越多分，越增多。\newline{}祂也用很多酒灌满水缸；他们汲出来，它虽多却缺乏：祂所赐的杯本身虽小，力量却极大，没有限量。祂是包含一切烈酒的杯，也是祂亲自在其中的奥迹！祂擘开的唯一之饼没有界限，祂调和的唯一之杯没有限量！\BibleRef{箴 9:5} 所播种的麦子，\BibleRef{若 12:24} 第三天发芽并充满生命的仓库。\BibleRef{玛 13:30} 属灵的饼，如同赐予者，属灵地使属灵的人活过来；而肉性地接受它的人，轻率接受，毫无益处。愿灵魂明辨地接受这恩宠之饼，作为生命的药物。\newline{}如果死祭是奉魔鬼的名献上的，\BibleRef{格前 10:20} 并且吃了，不是没有奥迹；对于祭献的圣物，我们更应当谨慎地施行这奥迹。那吃魔鬼名下的祭物的人，毫无异议地成为魔鬼的。那吃天上食粮的人，毫无疑问成为属天的！酒教导我们，因为它使亲近它的人变得像它：因为它很恨那喜爱它的人，使人醉而疯狂，并成为他的戏弄者！\BibleRef{箴 20:1} 光教导我们，因为它使眼睛——太阳的女儿——像它自己：眼睛借着光看见了赤裸，于是跑去羞怯地遮盖了那贞洁的人。\BibleRef{创 9:23} 至于那赤裸，是酒造成的，酒甚至对贞洁者也不表示怜悯！\par

首生者以欺骗者的武器披戴自己，为要以那杀人的武器，使人重获生命！祂以那杀害我们的树拯救了我们，以那使我们癫狂的酒，我们藉以变得贞洁！由亚当取出的肋骨，恶者以此夺去了亚当的心。从这肋骨中升起\BibleRef{创 3:15}一种隐藏的力量，斩断了撒殚如同达贡一般：因为在约柜中藏着一卷书，它呼喊并宣告那得胜者！于是有一个奥秘被启示出来，即达贡在自己的避难所中被推倒！\BibleRef{撒上 5:4}预象之后，应验来到：恶者在他所信赖之处被推倒！愿那来临者受赞颂，在祂身上，左与右的奥秘都得以成就。\BibleRef{玛 25:33}那在羔羊内的奥秘应验了，那在达贡内的预象也应验了。愿祂受赞颂，祂以真羔羊救赎了我们，并摧毁了我们的毁灭者，如同摧毁达贡一般！在十二月，当夜长之时，那无限之日为我们升起！在冬季，当世界昏暗时，那佳美者出来，欢愉世上众人！在冬季使大地荒芜时，童贞学会了生育。在十二月，使大地的劳作停息时，其中有了童贞的劳作。早期的羔羊，牧人从未在见到牧人之前见过它：至于真羔羊，在祂诞生的时节，有关祂的喜讯也急忙传给了牧人们。那老狼看见吃奶的羔羊，虽然祂隐藏了自己，狼却在他面前战栗；因为狼披上了羊皮，万有的牧者成了羊群中的羔羊，为使那贪婪者胆敢攻击温良者时，大能者能撕裂那吞噬者。\BibleRef{民 14:6}圣者肉身形体地居于胎中，又神性形上地居于心灵中。怀孕祂的玛利亚憎恶婚床：愿那灵魂不要犯奸淫，祂居住其中。因为玛利亚领悟了祂，便离开了她的未婚夫：祂居住在贞洁的童女中，如果她们领悟祂。\BibleRef{玛 5:28}聋子听不见巨雷，鲁莽的人也听不见诫命的声音。因为聋子在霹雳时惊惶，鲁莽的人在训诲的声音中也惊惶；如果可怖的雷声恐吓聋子，那么可畏的忿怒也会激动不洁的人！聋子听不见并非他的过错，但谁践踏（诫命）就是鲁莽。不时有雷声，但法律的声音每日如雷轰响。我们不要闭上耳朵，当它们的开口——作为敞开而非闭合以对抗它——控告我们时；听觉之门自然敞开，它要责备我们违抗自己意志的鲁莽。声音之门和口舌之门，我们的意志可以开启或关闭。让我们看看良善者赐予了我们什么；让我们聆听那大能的声音，不要闭上我们耳朵的门。\par

\OriginalTerm{声音}{Voice}成了身体，\OriginalTerm{至高者}{High One}的圣言成了血肉，愿光荣归于祂！耳朵啊，也请听祂；眼睛啊，也请看祂；手啊，也请触摸祂；口啊，也请享用祂！肢体与官能，你们应当赞美祂，祂前来使整个身体复生！玛利亚怀了那沉默的婴孩，而祂内却蕴藏着万种语言！若瑟怀抱着祂，而祂内蕴藏着比一切古老事物更久远的本性！\OriginalTerm{至高者}{High One}成为一个小婴孩，而祂内蕴藏着充足万有的智慧宝库！祂虽至高，却吮吸玛利亚的乳汁，而万物因祂的良善汲取生命！祂是\OriginalTerm{生命之乳}{Breast of Life}，也是生命的气息；死者从祂的生命中汲取而复活。没有空气的气息，无人能活；没有\OriginalTerm{圣子之能}{Might of the Son}，无人能存。祂赐生命的气息使万物复生，天上的与地下的诸灵都依赖它。当祂吮吸玛利亚的乳汁时，祂正以生命滋养万有。当祂躺在母亲怀中时，万物都躺在祂怀中。祂如婴孩般沉默，却使万物遵行祂的一切命令。因为除非藉着\OriginalTerm{首生者}{First-born}，无人能接近那\OriginalTerm{本体}{Essence}，祂与本体同等。祂在地上三十载，掌管万物，收纳天上与地下一切赞颂之祭。祂全然在深渊，又全然在至高之处！祂全然与万物同在，又全然与每个个体同在。当祂的身体在胎中成形时，祂的能力正塑造所有肢体！当圣子的成胎在母腹中形成时，祂自己正在母腹中塑造婴孩。然而，祂的身体在胎中软弱，祂的能力在胎中并非软弱！同样，祂的身体在十字架上软弱，祂的能力在十字架上也并非软弱。因为当祂在十字架上使死者复活时，祂的身体使他们复活，是的，更是祂的\OriginalTerm{旨意}{Will}；正如当祂全然居于胎中时，祂隐藏的旨意正巡视万有！请看，当祂全然悬挂在十字架上时，祂的能力仍使万物运转！因为祂使太阳昏暗，使大地震动；祂撕裂坟墓，使死者出来！请看，当祂全然在十字架上时，祂又全然无处不在！因此，祂全然在胎中，同时又全然在万有之中！当祂在十字架上使死者复活时，同样，当祂是婴孩时，祂正塑造婴孩。当祂被杀害时，祂开启了坟墓；\BibleRef{玛 27:52}当祂在胎中时，祂开启了胎胞。来吧，我的弟兄们，请聆听关于\OriginalTerm{奥秘者之子}{Son of the Secret One}的事，祂在身体中启示，而祂的能力却隐藏着！因为圣子的能力是自由的能力；胎胞不能束缚它，如同束缚身体一般！因为当祂的能力居于胎中时，祂正在胎中塑造婴儿！祂的能力环绕着那环绕祂的胎胞。因为若祂收回祂的能力，万物就会倾倒；祂的能力支撑万物；当祂在胎中时，祂并未松开对万有的掌握。祂亲身在胎中塑造了一个形象，并正在所有胎中塑造一切容颜。当祂在穷人中间渐渐长大时，祂从丰富的宝库滋养万有！当那膏抹祂的妇人膏抹祂时，祂正以自己的露和雨膏抹万有！贤士们带来没药与黄金，而祂内蕴藏着富足的宝库。贤士们带来的没药与香料，原是祂自己预备并创造的，他们将祂自己的东西献给了祂。正是出于祂的能力，玛利亚才能将托住万有的那一位抱在怀中！玛利亚所给祂的一切，都是从万物的巨大仓库中取来的！\BibleRef{耶 31:22}她把乳汁给祂，这乳汁原是祂自己预备的；她把食物给祂，这食物原是祂自己创造的！祂作为天主赐给玛利亚乳汁，又作为人子从她那里吮吸。她以双手怀抱祂，因为祂已虚己；她的手臂拥抱祂，因为祂已使自己微小。祂的威严谁能测量？祂使自己的度量缩小，成为一件衣裳。她为祂编织并穿上衣裳，因为祂已脱下自己的荣耀。她量度祂并为祂编织，因为祂已使自己微小。\par

海水载祢时，曾静止平息，若瑟的怀抱又如何能承受祢？地狱的腹怀怀了祢，却破裂了，玛利亚的子宫又如何能容纳祢？坟上的巨石被祢的大能冲破，玛利亚的手臂又如何能容下祢？祢降至卑微，为将众人复活！愿赞美从所有因祢而活者归于祢！谁能述说那隐藏者的圣子，祂降下，在母胎中穿上身体？祂出来，如婴孩吮吸乳汁，在孩童中，万有之主匍匐爬行。人们见祂在街上如幼童，而祂内住着众爱。孩童们明见围绕祂在街上；天使们暗自在敬畏中围绕祂。祂与幼童同乐如孩童，祂对天使可畏如统帅：祂对若翰可畏，使他解祂的鞋带：祂对罪人温柔，他们亲吻祂的脚！天使们以天使之眼看见祂；每人按自己的知识程度看见祂：每人按自己的辨识程度，如此认识那超乎万有者。圣父和祂自己乃是知识的全量，得以认识祂的本体！因为一切受造物，无论上界下界，都各得其所量知识；祂，万有之主，将一切赐给我们。祂使万有富足，却向万有索取利息。祂赐予一切，似无所缺，却又如匮乏者向万有索求。祂作为造物主赐我们牛群羊群，却又似缺乏而要求祭献。祂作为创造者使水变酒，却以穷人身份饮酒。在婚宴中，祂用己物调和了酒，祂调和了自己的酒，作为宾客请人喝。因祂的爱，祂使年迈的西默盎多活了日子，使他，一个凡人，能献上那赐生命者。靠祂的力量，西默盎抱起了祂；献上祂的那人，反被祂献给天主。祂在山上给梅瑟覆手\BibleRef{出 33:22}，又在约旦河从若翰领受了覆手。凭祂恩赐的力量，若翰虽属地，却能给属天者授洗。凭祂的力量，大地支撑了祂：大地几近溶解，祂的力量坚固了它。玛尔大给祂食物：她将祂所造的菜肴摆在祂面前。凡奉献者，都用祂自己的东西献了愿：他们将祂宝库中的东西摆在祂的桌上。\par

% structure: p0050 source=h2 target=section reason=relative-heading-rank; base=h2

\section{第四首赞歌}

这是带来诸般喜乐的月份；它是奴隶的自由，自由人的骄傲，城门的冠冕，身体的安宁，它也在其爱中把紫袍加于我们身上，如同君王。\par

这是带来诸般胜利的月份；它释放精神，制服肉体，在凡人中产生生命；它因爱使天主性寓居于人性中。\par

在此日，主以谦卑将荣耀换成羞辱；因为亚当作为叛逆者，将真理换成不义；善者怜悯了他，使那背逆者成义并改正。\par

愿每人驱除疲倦，因那位至尊并未厌倦为我们在胎中九月，以及在索多玛的狂人中三十年。\par

因为善者看见人类贫穷卑微，祂便设立节庆如宝库，向懒惰者敞开，使节庆激励懒惰者起来致富。\par

看！首生者为我们开启了祂的节庆如同宝库。一年中仅此一日独自开启那宝库：来，让我们获取，让我们从中富足，趁它们关闭之前。\par

警醒者有福了，他们已用强力夺取\BibleRef{玛 11:12}了生命的掠物。那是极大的耻辱，当人看见邻人取走并运出宝藏，而自己却坐在宝库中打瞌睡，以致空手而出。\par

在此节庆，愿我们每人加冕自己心灵的门户。圣神渴望这门户，好能进入居住，圣化它，祂遍巡所有门户，要看能从何处进入。\par

在此节庆，门户为门户而喜乐，圣者在圣殿中喜乐，声音在孩童口中回响，基督在其节庆中如壮士喜乐。\par

在圣子诞生时，王为税银登记所有人，使他们成为祂的债务人：那王来到我们中间，涂抹了我们的债卷\BibleRef{哥 2:14}，又以祂自己的名字另写一张，使祂成为我们的债务人。太阳延长了光明，以它所上升的度数预表了奥秘。自它上升已有十二日，今日是第十三日：这正是圣子诞生\BibleRef{出 12:3}和祂十二宗徒的准确预像。\par

梅瑟在尼散月十日封闭了一只羔羊；这是圣子进入母胎并于第十日将自己封闭其中的预像。祂在这太阳延长光明的月份从母胎出来。\par

黑暗被战胜，为宣告撒殚被战胜；太阳延长光明，为凯旋，因为首生者得胜。黑暗与黑暗者一同被战胜，我们的光明与更亮的光一同得胜！\par

若瑟爱抚圣子如同婴孩；他事奉祂如同天主。他因祂是善者而喜乐，又因祂是义者而畏惧，心中大感困惑。\par

\Quote{谁将至高者之子赐给我为子？我曾嫉妒祢的母亲，想要休她，却不知她腹中藏着巨大的宝藏，要突然使我这穷困者致富。达味王从我族中生出，戴上了王冠；而我却落得极其卑微，本为王，却成了木匠。然而王冠临到了我，因我怀中是万冠之冠！}\par

\Quote{谁赐给我这不孕者，使我怀了并生了这一位，祂是多样的一位，微小却伟大；因祂完全与我同在，又完全无处不在？}\par

\Quote{加俾额尔进入我这卑微之境的那日，他使我忽然从婢女变为自由人：因我本是祢神性的婢女，又成为祢人性的母亲，主与圣子啊！}\par

忽然之间，婢女因祢——君王之子——而成了君王之女。看，达味家中最卑微者，因祢——达味之子——而成为属天者所达至的天上之人！\par

我何等惊异，在我面前竟有一位孩童，比万有更年长！祂的眼目恒常仰望天上。祂口中咿呀之声，在我看来，显示其沉默在天主面前发言。\par

谁曾见过一位孩童，其整个存在洞悉一切？祂的视线如同那统管天上地下万物的主宰！祂的容貌如同那指挥万有的统帅。\par

我如何向祢——泉源——开启乳汁之泉？或如何向那从祢筵席滋养万有的祢提供滋养？我如何将襁褓带给那以荣光之射线为衣的祂？\par

我的口不知如何称呼祢——永生者的孩童：因若大胆称祢为若瑟之子，我战栗，因祢非他后裔；我又惧怕否认那与我订婚者之名。\par

祢既是一位的儿子，我岂可称祢为多人之子？因万千名号不足以称呼祢，祢既是天主子，又是人子，是的，达味之子，玛利亚之主。\par

谁使众口之主成为无口者？因我对祢的纯洁受孕，恶人毁谤了我。圣者啊，为祢的母亲发言吧！显个奇迹，好使他们确信我由谁怀了祢！\par

因祢的缘故，我也被憎恨，祢——万有之爱者。看，我虽怀了并生了那人类避难之所，却受迫害。亚当将喜乐，因祢是乐园之钥。\par

看，海洋向祢的母亲发怒，如同向约纳。看，黑落德那狂怒之浪，欲淹没诸海之主。我往哪里逃，祢必教导我——祢母亲的主。\par

与祢一同逃亡，好使我在祢内获得每处之生命。监牢与祢同在便非监牢，因在祢内人升到天上；坟墓与祢同在便非坟墓，因祢是复活！\BibleRef{若 11:25}\par

一颗非本性的光星，突然照耀；小于太阳又大于太阳，小于其可见之光，却大于其隐藏之力，因其奥秘。\par

晨星在黑暗中投出光明，引领他们如盲者，他们前来领受大光：献上礼物，获得生命，崇拜后返回。\par

在高天与深处，有两位宣讲者为圣子作证：明星在高处呼喊；若翰在下宣讲，两位宣讲者，一属天，一属地。\par

那在上者显示祂的本性来自尊威，那在下者也显示祂的本性来自人类。啊，伟大的奇迹，祂的天主性与人性各自被他们宣讲。\par

凡以为祂属地者，明星说服祂属天；凡以为祂属神者，若翰说服祂也有形体。\par

西默盎在圣殿中怀抱祂，轻摇祂，\Quote{慈悲者啊，祢来了，怜悯我的老年，使我的骨头安然进入坟墓。在祢内，我将从坟墓复活进入乐园！}\par

亚纳拥抱祂，以口亲祂的唇，圣神便停留在她唇上。正如依撒意亚的口缄默时，那接近他唇的火炭开启了他的口；同样，亚纳因祂口中的圣神而炽热，轻摇祂，\Quote{国度之子，谦卑之子，那听见而缄默，看见而隐藏，知道而未知者，天主，人子，荣耀归于祢的名。}\par

不孕者也听闻，奔跑前来，带着她们的礼品；贤士们带着宝物前来，不孕者带着礼品前来。礼物与财富突然堆积在穷人家中。\par

不孕者呼喊，如见所未见：\Quote{有福者啊，谁赐我见此景象，得见祢的婴孩，天地因祂而满盈！祢的果子有福，使不结果子的葡萄树结出一串。}\par

匝加利亚前来，开口尊贵之口喊道：\Quote{君王在哪里？我为祂生了那在祂面前宣讲的声音。万福，君王之子，我们的司祭职也将交给祢！}\par

若翰随同父母前来朝拜圣子，祂将荣光照耀他的面容；他不再如母腹中那样激动！伟大的奇迹，此处他朝拜，彼处他跳跃。\par

黑落德，那卑劣的狐狸，徘徊如狮子，又蹲伏如狐狸，嚎叫，因听见那狮子——按圣经前来坐王位者——的吼声。狐狸听说狮子是幼崽，尚在哺乳，便磨利牙齿，趁祂年幼时潜伏欲吞食狮子，待祂长大，祂口中的气息将毁灭他。\par

整个受造界都成了祂的口，为祂呼喊。贤士们以礼品呼喊！不孕者以子女呼喊，光星在空中呼喊：看，君王之子！\par

天开了，水静了，鸽子荣耀了祂，圣父的声音——比雷声更响——即时说：这是我所爱的儿子。天使传报祂，孩童以和散那向祂欢呼。\par

这些天上地下的声音传扬祂并高声呼喊。熙雍的沉睡未被雷声震醒，反而被触怒，起来杀了祂，因祂唤醒了她。\par

% structure: p0092 source=h2 target=section reason=relative-heading-rank; base=h2

\section{第五首赞美诗}

圣子诞生时，白冷有大呼喊，因天使降下，在那里赞美。他们的声音如巨雷；在赞美声中，沉默者前来，向圣子赞美。\par

那婴孩有福，厄娃与亚当在祂内恢复青春！牧人们也带着羊群中最好的礼物前来：甜奶、洁净的肉、合宜的赞美！他们分别献上：将肉给若瑟，奶给玛利亚，赞美给圣子！他们带来并献上一只羔羊给逾越节羔羊，首生者给首生者，祭品给祭品，时间之羔羊给真理之羔羊。何等美善的景象，羔羊献给羔羊！\par

羔羊在初生子前被献祭时哀叫。它赞颂那来到以解放羊群与牛脱离祭献的羔羊：是的，那逾越节羔羊，祂传授并带来了圣子的逾越节。\par

牧人们带着牧杖前来朝拜祂。他们以平安问候祂，同时预言说：\Quote{平安，牧人之王！}梅瑟的杖赞颂你的杖，万有的牧者；因为梅瑟赞颂你，尽管他的羊已成了豺狼，他的羊群仿佛成了巨龙，他的羊成了有牙的野兽。在那可怕的旷野中，他的羊群狂怒起来，攻击了他。\par

牧人们因此赞颂你，因为你使豺狼与羔羊在羊栈内和好。\newline{}噢，婴孩，你比诺厄更年长，又比诺厄更年轻，在洪涛中使一切在方舟内和好！\par

你的父亲达味为了一只羔羊杀了一头狮子。你，达味之子，杀死了那无形的豺狼，它谋杀了单纯的亚当，那只在乐园中吃草、哀叫的羔羊。\par

听到那赞颂的声音，新娘们被感动而自洁，童女们守贞，甚至少女们也端庄起来；她们成群前来，朝拜了圣子。\par

达味城中的老妇们来到达味之女面前；她们感谢说：\Quote{我们的乡国有福，因她的街道被叶瑟的光芒照亮！今日达味的宝座因你，达味之子，得以建立。}\par

老人们呼喊说：\Quote{那圣子有福，祂使亚当恢复青春；祂看见亚当衰老疲乏，那杀死他的蛇已蜕皮逃脱，便感到忧伤。那婴孩有福，在祂内亚当和厄娃都恢复了青春。}\par

贞洁的妇女说：\Quote{噢，有福的果实，请祝福我们腹中的果实；愿他们作为首生子奉献给你。}她们热情洋溢，预言她们的孩子，当他们为祂被杀时，就像初熟之果被采下。\par

不孕的妇女也爱抚祂，怀抱祂；她们欢喜说：\Quote{有福的果实，无婚姻而生，请祝福我们这些已婚者的子宫；怜悯我们的不孕吧，你这童贞的奇妙孩子！}\par

% structure: p0104 source=h2 target=section reason=relative-heading-rank; base=h2

\section{第六首赞歌}

那带着重担而来的使者有福，伟大的平安！圣父的心肠将祂带下给我们；祂没有将我们的债务带上去给祂，而是用自己的财富满足了那尊威。\par

智慧者有福，祂使神性与人性和好并结合。一位从上而来，一位从下而来，祂将本性如同药物一样限制，且身为天主的肖像，成了人。\par

那嫉邪者看见亚当本是尘土，那受咒诅的蛇已吞吃了他，便将健康注入那无味之物，使他成为盐，以之蒙蔽那受咒诅的蛇的眼睛。\par

仁慈者有福，祂看见乐园前的武器，挡住了通往生命树的道路；便来取了可受苦的身体，以便用祂肋旁的门户，打开进入乐园的路。\par

仁慈者有福，祂不偏向严厉，却用智慧不战而胜；为给人树立榜样，使人们能用德行和智慧明辨地得胜。\par

你的羊群有福，因为你是它的门，你是它的牧杖。你是它的牧者，你是它的饮品，你是它的盐，是的，你是它的眷顾者。万福，\OriginalTerm{独生子}{Only-Begotten}，你丰盛地承载了各样安慰！\par

农夫们前来朝拜生命的农夫。他们欢喜地向祂预言说：\Quote{农夫有福，祂耕耘心中的田地，将祂的麦子收进生命的仓里。}\par

农夫们前来荣耀那从叶瑟的根和干发出的葡萄园，荣耀葡萄的童贞果实。\Quote{愿我们成为你新酒的器皿，这酒更新一切。}\par

\Quote{愿在我爱人的葡萄园——它曾结出野葡萄——在你内得享平安！将你的枝条接在它的葡萄树上，让它完全承载你的祝福，结出能平息那葡萄园主人愤怒的果实，因祂正威胁它。}\par

因若瑟的缘故，工匠们来到若瑟之子面前说：\Quote{你的诞生有福，你是工匠之首，方舟曾带着你的印记，在此之后，会幕也按它建造，只是暂时的。}\par

\Quote{我们的手艺赞颂你，你是我们的荣耀。请造出轻省的、是的，容易的轭，给背负它的人；请造出没有虚伪的、充满真理的量器；是的，请按正义造出量器；这样，卑污之人可以因此被定罪，完全之人可以因此被宣告无罪。请用它们称量仁慈与真理，公义者，如同法官。}\par

\Quote{新郎与新妇一同欢喜。那婴孩有福，祂的母亲是圣神的新娘！那婚宴有福，你在其中，当酒突然短缺时，在你内它又丰盈起来！}\par

孩子们呼喊说：\Quote{祂有福，祂成为我们的兄弟，在街市中与我们同行。那日子有福，它借着树枝荣耀生命树，祂使自己的尊威降卑，降至我们童年的年岁！}\par

妇女们听说一位童贞女要怀孕生子：尊贵的妇女们希望你会从她们中兴起；是的，高贵的女士们希望你从她们中生出！你的尊威有福，它谦卑自己，从贫贱中兴起！\par

是的，怀抱祂的少女们预言说：\Quote{无论我是被憎恨还是美丽，或是卑微，在你面前我毫无玷污。我已为生产之床照顾你。}\par

撒辣曾安抚依撒格，他如奴隶\BibleRef{创 22:6}，在肩上背着君王他主人的肖像，即祂十字架的记号；是的，他的手上缠着绷带并承受痛苦，那是钉痕的预像。\par

辣黑耳向她的丈夫哭求说：\Quote{给我儿子。}\BibleRef{创 30:1}玛利亚有福，在她腹中，你虽未被求，却圣洁地居住，噢，恩赐，你倾注在领受者身上。\par

亚纳以苦泪祈求孩子；\BibleRef{撒上 1:7}撒辣与黎贝加以誓言和话语，依撒伯尔也以祈祷，她们长期苦恼之后，如此得到了安慰。\par

玛利亚是可赞美的，她未发愿、未祈祷，以童贞之身怀孕并生下了主，即她所有同伴之子中的主，这些同伴或曾或将要成为贞洁而正义的人、司祭和君王。\par

谁曾像玛利亚那样在怀中抚育儿子？谁曾敢称她的儿子为造物主之子、创造者之子、至高者之子？\par

谁敢在祈祷中向她儿子说话？\Quote{哦，祢作为天主是祢母亲的信赖，作为人则是她的爱子和儿子，恐惧与爱情中，祢的母亲理应站立在祢面前！}\par

% structure: p0126 source=h2 target=section reason=relative-heading-rank; base=h2

\section{颂歌第七首}

造物主之子如同祂的父一样是造物主！祂为自己造了一个纯净的身体，给自己穿上，然后出来，用祂从父那里慈悲带来的荣耀，遮盖了我们的软弱。\par

从默基瑟德大司祭那里，有一束牛膝草来到祢面前；从达味家族那里，有宝座和王冠；从亚巴郎那里，有族系和家室。\par

求祢因自己的缘故，成为我的港湾，哦，浩瀚的海洋！看，祢的父达味的圣咏，以及先知们的话语，都如同船只般来到我这里。\par

祢的父达味在第一一零篇圣咏中，将两个数字如同冠冕般编织给祢，来到祢面前，得胜者！祢将戴上这些冠冕，登上宝座，坐在其上。\par

一个伟大的冠冕是编织在百数中的，其中祢的天主性被加冕；一个小小的冠冕是十数的那个，它加冕祢人性的头，胜利者！\par

为了祢的缘故，妇女们追求男子。塔玛尔渴望那鳏居者，卢德爱上了一位老人，是的，那曾俘虏人的辣哈布也被祢所俘获。\par

塔玛尔出去了，在黑暗中\EditorialNote{创世纪38章}偷窃了光明，在不洁中偷窃了圣者，并因暴露自己的裸体而进入，偷窃了祢，荣耀者，使纯洁者从污秽者中诞生。\par

撒殚看见她，就战栗，急忙去扰乱她。他将审判放在她心中，她却不害怕；石刑和刀剑，她也不战栗。那教导奸淫者阻止了奸淫，因为他本是祢的阻碍者。\par

因为塔玛尔的奸淫是圣洁的，为祢的缘故。祢是她所渴求的，纯洁的源泉。犹大欺骗了她，不让她饮祢。干渴的子宫从祢的源头偷窃了祢的露珠。\par

她为祢成了寡妇，她渴望祢；她急忙也为祢成了娼妓，她热切地想要祢，并因她爱的是祢而成圣。\par

愿塔玛尔欢欣，因为她的主来了，并使她的名声因她奸淫的儿子而显扬！确实，她给他起的名字\BibleRef{创 38:29}是在呼唤祢到她这里来。\par

为了祢，尊贵的妇女们自取其辱，祢将贞洁赐给众人！她在路中将祢偷走，祢开辟了通往天国的道路！因为她偷的是生命，刀剑不能杀死她。\par

卢德为了祢的缘故，在打谷场上躺在一个男人旁边；她的爱使她勇敢，为了祢，哦，祢教导所有忏悔者勇敢。她的耳朵拒绝听任何声音，为了听祢的声音。\par

那发光的火炭升入波阿次的床榻，躺卧在那里，看见了那位大司祭，在他的腰中隐藏着为他香料而燃烧的火！她急忙成为波阿次的母牛，要生下祢，肥美的牛犊。\par

她因爱祢而去拾穗，收集麦秆。祢迅速以她的谦卑赏赐了她；代替麦穗，祢使君王之根发出；代替麦秆，祢使生命之穗从她而生。\par

% structure: p0142 source=h2 target=section reason=relative-heading-rank; base=h2

\section{颂歌第八首}

为使祢的复活在不信者中被相信，他们将祢封在坟墓里，并设了守卫；因为他们确实为祢封了坟墓并设了守卫，永生之子！\par

当他们埋葬祢后，如果他们将祢弃置不顾，离去，就有机会撒谎说他们偷了，哦，众生的复苏者！当他们巧妙地封了祢的坟墓，他们反而使祢的荣耀更大。\par

因此，达尼尔和拉匝禄是祢的预像：一个在穴中被外邦人封住，一个在坟墓中被百姓打开。看，他们的标记和封印谴责了他们。\par

如果他们让祢的坟墓敞开，他们的嘴本会是张开的。但他们离开了，因为他们关闭并封住了祢的坟墓，也封住了自己的口。是的，他们关闭了，当他们愚蠢地遮盖了祢的坟墓时，所有诽谤者都遮盖了自己的头。\par

但在祢的复活中，祢说服他们关于祢的诞生；因为子宫被封住，坟墓被关闭：在子宫中同样纯洁，在坟墓中同样活着。子宫和坟墓都被封住，它们为祢作证。\par

母腹和阴间大声呼喊祢的诞生和复活：被封住的子宫怀了祢，关闭的阴间交出了祢。无论是子宫怀祢，还是阴间交出祢，都不是按着本性！\par

坟墓被封住，他们将祢托付给它，让它保守死者；子宫是童贞的，无人认识。童贞的子宫和封住的坟墓，如同号角，向聋耳的人民宣讲祂。\par

封住的母腹和关闭的磐石都在控诉者中。因为他们毁谤受孕是出于人的种子\OriginalTerm{人的种子}{seed of man}，复活是出于人的盗窃；封印和印鉴定他们的罪，并辩称祢是从天上来的。\par

人民站在祢的诞生和祢的复活之间。他们毁谤祢的诞生，祢的死亡定了他们的罪；他们否认祢的复活，祢的诞生驳斥了他们；就像两个摔跤手阻止了诽谤的嘴。\par

至于厄里亚，他们去搜山：\BibleRef{列下 2:16} 当他们在地上寻找他时，他们更加确认他被接去了。他们的搜寻证明了他被接去，因为他们没有找到他。\par

因此，如果预先得到厄里亚升天警告的先知们尚且好像怀疑他的上升，那么不洁之人对圣子的诽谤该有多么厉害！祂通过他们自己的守卫使他们相信祂已复活。\par

主啊，没有人知道该给祢的母亲取什么名字。若称她为童贞，她的孩子就在那里；若称她为已婚，却无人认识她！如果没有人能理解祢的母亲，谁又能理解祢呢？\par

因为她独然是祢的母亲；与众人一起，是祢的姊妹。她是祢的母亲，她是祢的姊妹。她与贞洁的妇女一同是祢的聘妻。祢以万有装饰了她，祢，祢母亲的装饰。\par

因为她在祢来临之前，按本性是祢的新娘；在祢来临之后，她不是按本性怀孕祢，\OriginalTerm{圣者}{Holy One}啊，当她圣洁地生下祢时，仍是\OriginalTerm{童贞}{Virgin}。\par

\Person{玛利亚}在祢身上，主啊，获得了所有已婚妇女的尊荣。她未结婚就在腹中怀了祢。她乳房中有乳汁，不是按自然的方式。祢使干渴之地突然成为乳泉。\par

若她怀抱祢，祢威严的目光使她的负担轻省；若她给祢吃，是因祢饿了；若她给祢喝，是因祢渴了；若她甘心拥抱祢，祢，\OriginalTerm{慈悲之炭}{coal of mercies}，保全了她的胸怀。\par

\Quote{\OriginalTerm{主}{Lord}进入了她，成了仆人}：\Quote{\OriginalTerm{圣言}{Word}进入了她，在她内成为静默}；\Quote{雷声进入了她，祂的声音静止了}：\Quote{\OriginalTerm{万有的牧者}{Shepherd of all}进入了她，祂在她内成为\OriginalTerm{羔羊}{Lamb}，咩咩叫着出来。}\par

祢母亲的腹中改变了万物的秩序，祢是安排万物的！富者进入，祂贫穷地出来；至高者进入，祂卑微地出来。\OriginalTerm{光明}{Brightness}进入了她，穿上衣服，出来时成了\OriginalTerm{卑贱的形像}{despised form}。\par

\OriginalTerm{全能者}{The Mighty}进入，从腹中就穿上了恐惧。那赐食物给万有的进入，得了饥饿。那赐饮给万有的进入，得了干渴。\OriginalTerm{万有的衣被者}{Clother of all}赤身露体地从她出来。\par

那些在\WorkTitle{\OriginalTerm{耶肋米亚哀歌}{Lamentations of Jeremiah}}中哭泣的\OriginalTerm{希伯来妇女}{daughters of the Hebrews}，不用她们圣经中的哀歌，却用她们自己书中的催眠曲；她们话语中隐藏的能力在预言着。\par

\Person{厄娃}从\OriginalTerm{阴府}{Sheol}举目，在那一天欢喜，因为\OriginalTerm{她女儿的儿子}{Son of her daughter}作为\OriginalTerm{生命之药}{medicine of life}降下，为复活\OriginalTerm{祂母亲的母亲}{mother of His mother}。\OriginalTerm{受赞美的婴孩}{Blessed Babe}，祢击碎了那击伤她的\OriginalTerm{蛇}{Serpent}的头！\par

她从俊美的\Person{依撒格}的青年岁月中看见了祢的预像。因为\Person{撒辣}看见祢的预像停留在他童年，就称他，说：我誓愿的孩子，在你内隐藏着誓愿之主。\par

\OriginalTerm{纳齐尔人}{Nazarite}\Person{三松}预表了祢工作的预像。他撕裂了狮子，那死亡的形像，祢毁灭了它，并从它的苦中为人类引出生命的甘甜。\par

\Person{亚纳}也拥抱了\Person{撒慕尔}；因为祢的正义隐藏在他内，他将\Person{阿加格}剁成碎块，作为恶者的预像。他为\Person{撒乌耳}哭泣，因为祢的良善也预表在他身上。\BibleRef{撒上 2:26}\par

祢何其温良！祢何其刚强，\OriginalTerm{婴孩}{Child}啊！\BibleRef{路 2:52}祢的审判是强大的，祢的爱是甘甜的！谁能抵挡祢？祢的圣父在天上，祢的母亲在地上；谁能述说祢？\BibleRef{依 53:8}\par

若有人寻求祢的\OriginalTerm{本性}{Nature}，它隐藏在天上\OriginalTerm{天主性}{Godhead}强大的怀中；若有人寻求祢\OriginalTerm{可见的身体}{visible Body}，它安放在他们眼前\OriginalTerm{玛利亚卑微的怀中}{lowly bosom of Mary}。\par

心神在祢的世代间游荡，\OriginalTerm{富有的那一位}{Rich One}啊！厚重的皱褶覆盖着祢的天主性。谁能测透祢的深处，祢这\OriginalTerm{大海}{great Sea}，使自己变小！\par

我们来见祢为\OriginalTerm{天主}{God}，看！祢却是\OriginalTerm{人}{man}；我们来见祢为人，\OriginalTerm{祢天主性的光明}{Light of Your Godhead}就照耀出来！\par

谁能相信祢是\OriginalTerm{达味王位的继承者}{Heir of David's Throne}？从祂所有的床铺中，祢继承了一个\OriginalTerm{马槽}{manger}；从祂所有的宫殿中，一个\OriginalTerm{洞穴}{cave}降临到祢。代替祂的战车，或许一匹普通的驴驹降临到祢。\par

祢何等无畏，婴孩啊，祢让所有人都得到祢来怀抱：祢向每个遇见祢的人微笑；每个看见祢的人都因祢而喜乐！祢的爱如同一个渴慕人类的人。\par

祢不区分祢的父辈与陌生人，也不区分祢的母亲与婢女，也不区分喂养祢的人与不洁的人。这是祢的鲁莽还是祢的爱，爱万有的那一位啊？\par

什么促使祢让所有看见祢的人都得到祢，无论贫富？祢帮助了那些没有呼求祢的人。祢如此渴慕人类，这是从何而来？\par

祢的爱何其伟大：若有人责备祢，祢不恼怒；若有人恐吓祢，祢不畏惧；若有人对祢发出嘘声，祢不觉烦恼！祢超越了报复伤害者的律法。\par

\Person{梅瑟}是温良的，但他的热诚是严厉的，因为他争战并杀戮。\Person{厄里叟}也使一个孩子复活，却用熊撕裂了许多孩子。你是谁，孩啊？你的爱比众先知的爱更伟大？\par

\OriginalTerm{哈加尔的儿子}{son of Hagar}踢了\Person{依撒格}。他忍受了，默不作声，他的母亲却嫉妒。祢是他的奥秘，还是他是祢的预象？祢像\Person{依撒格}，还是他像祢？\par

% structure: p0178 source=h2 target=section reason=relative-heading-rank; base=h2

\section{第九首}

来休息，在祢母亲的怀中安静，\OriginalTerm{光荣者的儿子}{Son of the Glorious}。鲁莽不适合君王之子。\OriginalTerm{达味之子}{Son of David}啊，祢是光荣的，却又是\OriginalTerm{玛利亚之子}{Son of Mary}，祢将祢的美貌隐藏在内室中。\par

祢像谁？喜乐的婴孩，美丽的小家伙，祢的母亲是童贞，祢的圣父是隐藏的，连\OriginalTerm{色辣芬}{Seraphim}都不能注视祢？告诉我们祢像谁，\OriginalTerm{恩慈者的儿子}{Son of the Gracious}！\par

当愤怒的人来看祢，祢使他们喜乐；他们彼此交换微笑：愤怒的人在祢内变得温和，甘甜的那一位啊。祢是有福的，小家伙，因为在祢内连苦的都变甜了。\par

谁曾见过一个婴孩，当在怀抱中对靠近他的人喜乐，看！祂自己却伸手给远处的人？美丽的景象：一个孩子，为每个人思虑，好让他们看见祂！\par

有烦恼的人来见了祢，他的烦恼就飞走了。有焦虑的人，因祢忘了他焦虑；饥饿的人因祢忘了食物；有差事的人，因祢迷失了，忘了他的行程！\par

安静吧，让众人去做他们的工作！祢是穷人的儿子，从祢自己学习：所有穷人都不得不放下工作来。爱人的那一位，祢用祢的喜乐将人们联结在一起。\par

\Person{达味}，那位威严的君王，拿起树枝，在节庆中在孩童间跳舞，献上赞美。这不是祢的父亲\Person{达味}的爱在祢里面温暖吗？\par

那个\OriginalTerm{撒乌耳的女儿}{daughter of Saul}！她父亲的魔鬼在她内说话：她称那位威严的君王为卑劣之人，因为他给了她民间的长老一个榜样，在赞美祢的日子，与孩子们一起拿起树枝。\par

谁敢将急躁的罪名加于祢？看，撒乌耳的女儿因嘲笑约柜而断绝了生育；她的口出言讥讽，口舌的报应便是不孕。\BibleRef{撒下 6:23}\par

让口舌因亵渎而震颤，免得它们被封闭！熙雍女子，向祂闭口吧，因为祂是达味之子，在你们面前欢跃。不要像撒乌耳的女儿那样，她的后裔已经断绝。\par

因为厄里亚克制了身体的欲望，便使雨停止降于那淫乱者；因为他制服了他的身体，便使露水远离那纵欲者，他们的泉源肆意涌流。\par

因为肉身情欲的隐秘之火没有在他内掌权，上天的火就服从了他。既然他在地上制服了肉欲，便升到那圣洁安息之处。\par

厄里叟同样克制了自己的身体，却使死人复活。死人复活是因着一种超乎常情的圣化而按常理发生；他复活了孩子，因为他使自己灵魂纯洁如断奶的孩子。\par

梅瑟断绝并与妻子分离，就在那淫妇面前分开了大海。漆颇辣虽为外教司祭之女，却持守圣洁；亚巴郎的女儿却与牛犊行淫。\par

% structure: p0193 source=h2 target=section reason=relative-heading-rank; base=h2

\section{颂歌第十首}

我当从祢开始发言，祢是元首，开始了万物。\BibleRef{默 3:14} 我，是的，我将开口，但祢充满我的口。我向祢是大地，祢是农夫。将祢的声音撒在我身上，\BibleRef{希 6:7} 祢曾将自己撒在祢母亲胎中。\par

希伯来所有贞洁的女子，首领们的处女女儿，都因我而惊讶！穷人的女儿因祢而招嫉，弱者的女儿因祢而招妒。谁将祢赐给了我？\par

\Quote{富者的儿子，厌弃富妇胸怀的那位，谁引导祢到穷人那里？因为若瑟贫穷，我也缺乏，然而祢的商人来了，带着黄金，到了穷人的家里。}\par

她看见了贤士们：因他们的奉献，她的歌声加增了。看！祢的朝拜者围绕着我，祢的奉献环抱着我。愿婴儿受赞颂，祂使祂母亲成为祂言语的竖琴。\par

正如竖琴等待主人，我的口等待祢。愿祢母亲的舌头说出令祢喜悦的话。既然我从祢学会了新的怀孕，愿我的口在祢内学会，新生儿子，一首新的赞歌。\par

如果阻碍对祢不算阻碍，因为困难在祢是容易的，正如无婚的胎怀了祢，无精的腹生了祢，小小的口也易于倍增祢的伟大荣耀。\par

看！我受压迫、受轻蔑，却仍喜乐；我的耳朵充满了责斥和嘲讽。但对我来说，忍受这一切是小事，因为祢的一个安慰能驱散万千烦恼。\par

\Quote{既然祢不轻视我，儿子啊，我的面容就明亮。我因怀孕而受毁谤，却生出了真理，祂称我为义。如果塔玛尔因犹大而被称为义，我岂不更因祢而被称义？}\par

祢的父亲达味在祢未来之前已在诗篇中歌颂祢，说祢将得着舍巴的黄金。他所歌颂的这篇诗，看！当祢还是婴儿时，它已真的将没药和黄金堆在祢面前。\par

他所写的一百五十篇诗篇，都在祢内得了滋味，因为一切预言的话语都需要祢的甘甜；没有祢的盐，一切智慧都淡而无味。\BibleRef{约 6:6}\par

% structure: p0204 source=h2 target=section reason=relative-heading-rank; base=h2

\section{颂歌第十一首}

（童贞母亲对圣婴说。）\par

我不会嫉妒，我的儿子，因为祢与我同在，也与众人同在。向承认祢的人作天主，向服事祢的人作主，向爱祢的人作兄弟，为使祢赢得众人。\par

当祢在我内居住时，祢也在我之外居住；当我公开生下祢时，祢隐藏的能力并未离开我。祢在我内，又在我外，祢使祢的母亲惊异。\par

因为当我看见祢外在的形像在我眼前，那隐藏的形像便\Quote{在我思想中}被映照出来，圣者啊。在祢可见的形像中，我看见亚当；在祢不可见的形像中，我看见与祢结合的祢的父。\par

那么，祢只向我一人以两种形像显示祢的美吗？愿饼映出祢，也愿思想映出祢；也住在饼中，住在吃饼的人身上。愿祢的教会既秘密又公开地看见祢，如同祢的母亲。\par

憎恨祢的饼的人，如同憎恨祢的身体。远处渴望祢的饼的人与近处爱慕祢的肖像的人，是一样的。在饼和身体中，先来的和后来的都看见了祢。\par

然而，祢可见的饼远较祢的身体宝贵；因为不信的人也见过祢的身体，却没有见过祢永生的饼。远处的人却欢喜！他们的份完全轻看近处的人的份。\par

看！祢的肖像在饼中藉葡萄的血映出；它也藉着爱的指头、信德的色彩映在心上。愿那藉着真理的肖像使偶像消逝的那一位受赞美。\par

祢并非如此的人子，使我向祢唱普通的摇篮曲；因为祢的怀孕是新的，祢的诞生是奇妙的。没有圣神，谁能向祢歌唱？一种新的预言的火热在我内燃烧。\par

我怎能称祢为我们中的陌生人？我该称祢儿子？还是称祢兄弟？\BibleRef{玛 12:50} 我该称祢丈夫？还是称祢主？孩子，祢使祢的母亲从水里第二次诞生。\par

因为我是祢的姊妹，同属我们双方之父达味的家族。再者，因祢的怀孕，我是祢的母亲；因祢的圣化，我是祢的配偶；因祢用以赎买并洗礼我的血和水，我是祢的婢女和女儿。\par

至高者之子来住在我内，我成了祂的母亲；又如同第二次生产，我生了祂，祂也藉第二次生产而生了我，因为祂穿上祂母亲的衣服，她将自己的身体披上祂的荣耀。\par

塔玛尔，出自达味家族，被阿默农玷辱；贞洁从他们二人中失落了。我的珍珠没有失落：它被存放在祢的宝库中，因为祢已将之佩戴。\par

她丈夫兄弟的香气从塔玛尔那里逸出，因为塔玛尔偷了他的香料。至于若瑟的新娘，连他的气息也从她的衣裳散发出来，因为她孕育了肉桂。\BibleRef{歌 4:14} 你的怀孕于我是一道火墙，啊，圣洁的儿子。\par

小花萎谢了，因为荣光之百合的香气\BibleRef{歌 2:1}太大。香料宝库不缺花朵或其香气！肉身退避，因为它感知到胎中有一个来自圣神的怀孕。\par

女人伺候男人，因为他是她的头。若瑟起来伺候他的主，那主是在玛利亚内。司铎因你的圣洁而伺候你的约柜。\par

梅瑟携带了主所写的石版，若瑟则携带了那纯净的版块，创造者之子居住其中。石版止息了，因为世界充满了你的教训。\par

% structure: p0222 source=h2 target=section reason=relative-heading-rank; base=h2

\section{赞美诗第十二首}

我所怀的婴孩怀抱着我，玛利亚说，祂垂下双翼，将我拿起放在祂的翎羽之间，升入空中；并给了我一个许诺，说高天和深渊都将属于我子。\par

我曾看见加俾额尔称祂为主，也看见大司祭，那位年迈的仆人，怀抱祂、背负祂。我曾看见贤士们俯伏朝拜，也看见黑落德惊惶不安，因为君王已经来临。\par

撒殚也曾扼杀婴孩，好让梅瑟灭亡，\BibleRef{出 1:16}他又杀害婴孩，好让那永生者死去。祂逃往埃及，祂来到犹太地，为要劳苦奔波；他企图捕捉那将要捕捉自己的人。\par

厄娃在其童贞中穿上了羞耻的树叶：你的母亲在她的童贞中穿上了那足敷万有的荣光之衣。她将身体的小小衣裳给了那覆盖万有者。\par

那心中和意念中有你者，是有福的！因你，君王之子，她成了君王的宫殿；因你，大司祭，她成了至圣所！她没有家庭的烦扰，也没有丈夫的搅扰！\par

厄娃再次成了一个巢穴和一个洞穴，给那受诅咒的蛇，它进入并住在她内。它邪恶的计谋成了她的食物，使她变为尘土。你是我们的食粮，你也属于我们的种族，是我们的荣光之衣。\par

凡有圣洁者，若遇危险，看！这里有他的守护者！凡有邪恶者，看！这里有他的赦免者！凡有魔鬼者，这里有它的驱逐者！凡有疼痛者，看！这里有包扎他们伤口者。\par

凡有儿子者，让他来，成为我良人的弟兄！\BibleRef{玛 12:15}凡有女儿或本族少女者，让她来，成为我荣耀者的新娘！凡有仆人者，让他释放他，使他来服侍他的主。\par

负你轭的自由之子，我子，将得一份赏报；而负两个主人之轭的奴隶，即负上主和下人，他有两份祝福，两个轭有两份赏报。\par

自由的女子，我子，是你的婢女：若那为奴的也服事你，在你内她是自由的；在你内她将得安慰，因为她得了释放；隐藏的苹果积存在她怀中，\BibleRef{歌 2:3}若她爱你！\par

啊，贞洁的女子，切望我的良人，好让祂住在你内；你，那不洁的，也为要祂圣化你！你们各教会，也为要那创造者之子，祂来更新万有，好装饰你们！\par

祂接纳了那愚昧的，他们敬拜并服事众星；祂更新了那因亚当而陈旧的大地，亚当犯罪而衰老了。新的创造是其更新者的作品，那自足者修复了身体连同其意志。\par

来吧，瞎眼的，不花钱就得光明！来吧，瘸腿的，接受你们的脚！聋哑的，接受你们的声音！还有那断手的，来吧；残废的也要接受他的手。\par

是创造者之子，祂的宝库充满各种援助。凡没有眼珠的，来就祂，祂用泥造物又改变它，造肉身，开启眼睛。\par

凭一小块泥土，祂显明亚当是用祂的手形成的；死者的灵魂也为祂作证，人的气息是由祂吹入的；凭最后的见证，祂被证实为那最初者的儿子。\par

你们众癞病人，聚集前来，不劳苦就得洁净。因祂不会像厄里叟那样洗你们，他在河中浸了七次；祂也不会像司铎那样用洒水烦扰你们。外邦人和异乡人都已投奔大医师。\par

异乡人的等级在君王之子面前无立足之地；主并不对祂的仆人见外，也并未隐藏祂是万有之主。若义者使身体生癞，而祢洁净它；那么，身体的创造者憎恨身体，但祢爱它。\par

若它不是祢的塑造，作为义者，祢就不会医治它；\BibleRef{申 32:39}若它不是祢的受造物，健康时，祢就不会打击它。祢加在其上的惩罚，与祢医治的痛苦，都宣告祢是创造者之子。\par

% structure: p0241 source=h2 target=section reason=relative-heading-rank; base=h2

\section{赞美诗第十三首}

（参看赞美诗第二首。\Emph{主显节}。）\par

1. 在那被称为塞姆哈的君王时代，我们的主在希伯来人中兴起；塞姆哈与德恩哈统治且来临，君王在地上，圣子在天上；愿祂的统治受赞美！\par

2. 在那将人登入死亡册的君王时代，我们的救赎主降下，将人登入生命册。祂登记，他们也登记：在高天祂登记我们，在地上他们登记祂。荣耀归于祂的名！\par

3. 在那名为塞姆哈的君王时代，预像与现实相遇，君王与君王，塞姆哈与德恩哈。祂肩上的十字架，是祂国度的标志。愿背负它的祂受赞美！\par

4. 三十年间祂贫困地行走在地上！让我们编织一切节拍的赞美之声，弟兄们，为主之年，如同三十顶冠冕献给三十年。愿祂的降生受赞美！\par

5. 在第一年，那是宝库的统领和丰盛祝福的分施者，愿那曾在荣耀中背负圣子的革鲁宾与我们一同赞美祂！祂舍弃荣耀，劳苦找到了迷失的羊。愿感谢归于祂！\par

6. 在第二年，愿色辣芬与我们更甚地赞美祂。他们曾宣告圣子为圣，旋即看见祂在反对者中受辱骂；祂忍受了轻蔑，并教导了赞美。愿光荣归于祂！\par

7. 在第三年，愿弥额尔及其随从，就是那些在至高之处服事圣子的，与我们一同赞美祂。他们看见祂在地上服事，洗脚，洁净灵魂。愿祂的谦卑受赞美！\par

8. 在第四年，愿整个大地与我们一同赞美祂。对圣子而言它未免太小，但它感到惊奇，因为它看见自己在那极其简陋的床榻上接待了祂。祂充满了这床榻，也充满了天堂。愿威严归于祂！\par

9. 在第五年，太阳照耀了大地。愿它以它的气息赞美我们的太阳，祂将自己的广阔降低，并谦卑了自己的大能，使那不可见灵魂的锐利眼睛能够看到祂。愿祂的光明受赞美！\par

10. 在第六年，愿整个空气与我们一同赞美祂，在它的广阔空间中万物得享光荣；它看见自己伟大的主竟在一个小小的怀抱中成为一个小小的婴孩。愿祂的尊贵受赞美！\par

11. 在第七年，云和风与我们一同喜乐，将露珠洒在花朵上，因为它们看见圣子奴役了自己的光明，并忍受了耻辱和可憎的唾污。愿祂的救赎受赞美！\par

12. 在第八年，愿田野赞美祂，它们从祂的泉源哺育自己的果实。它们敬拜，因为它们看见圣子被抱在怀中，纯洁者吮吸纯洁的乳汁。愿祂的善意受赞美！\par

13. 在第九年，愿大地荣耀她造物主的全能，祂起初将种子撒在她内，使她能产生一切产物；因为它看见玛利亚——一块干渴的土地——结出了一个奇妙的、甚至可惊异的婴孩的果实。[然后]它更甚地赞美祂，因祂是一切美善的汪洋。愿祂受高举！\par

14. 在第十年，愿西乃山荣耀祂，它曾在它的主面前战栗。它看见他们拿起石头攻击它的主；祂接受了石头，祂却要在磐石上建立祂的教会。愿祂的建造受赞美！\par

15. 在第十一年，愿大海赞美圣子量度它的拳头\BibleRef{依 40:12}，它感到惊愕，看见祂下来，在少量的水中受洗，并洁净了受造物。愿祂高尚的行为受赞美！\par

16. 在第十二年，愿圣殿赞美祂，它曾看见那孩童坐在老人中间：当节日的羔羊在祂的庆节中咩叫时，司祭们都静默无声。愿祂的赎罪祭受赞美！\par

17. 在第十三年，愿冠冕与我们一同赞美那得胜的君王，祂死去，并被戴上荆棘冠冕，又在亚当头上系上一顶伟大的冠冕在祂右边。愿祂的宗徒职分受赞美！\par

18. 在第十四年，愿埃及的逾越节赞美那前来并跨越万有的逾越节，它代替法郎淹没了军团，代替马匹扼杀了魔鬼。愿祂的报复受赞美！\par

19. 在第十五年，愿贪食者的羔羊赞美祂：因为我们的主远非如梅瑟那样宰杀它，反而以祂自己的血救赎了人类。那养活万有的，为万有而死。愿祂的父受赞美！\par

20. 在第十六年，愿麦子以其预像赞美那位农夫\BibleRef{若 12:24}，祂将自己的身体撒在不毛之地，因为它覆盖一切，铺展自身，产出新饼。愿那纯洁者受赞美！\par

21. 在第十七年，愿葡萄树赞美那装饰它的主。祂栽种了葡萄园，灵魂如同葡萄苗。祂赐给葡萄园和平，却毁灭了那结出野葡萄的葡萄园。愿那拔除它者受赞美！\par

22. 在第十八年，愿那被林中野猪吃尽的葡萄树赞美那真正的葡萄树\BibleRef{若 18:9}，它修剪了自己，保存了自己的果实，并将果实带给葡萄园的主人。愿祂的收成受赞美！\par

23. 在第十九年，愿我们的面酵赞美那真正的面酵，它在那些错误者中间活动，将他们全部驱赶到一起，并以一个教义使他们同心。愿祢的教义受赞美！\par

24. 在第二十年，愿盐赞美祢的活体\BibleRef{谷 9:49}，藉此所有信友的身体和灵魂被盐腌，而信德是人的盐，藉此他们得以保存。愿祢的保存受赞美！\par

25. 在第二十一年，愿旷野的水赞美祢\BibleRef{出 15:25}。它们对远处的人是甜的，对近处的人是苦的，那些人不服事祂。在旷野中，[选民]和外邦人是苦的，祂毁灭了他们。他们被救赎他们的十字架变甜了。愿祢的甘饴受赞美！\par

26. 在第二十二年，愿武器和刀剑赞美祢：它们不足以杀死我们的仇敌。是祢杀了他，甚至祢安上了西满的剑所砍下的耳朵。愿祢的医治受赞美！\par

27. 在第二十三年，愿驴赞美祂，它把自己的驴驹给了祂骑，解开了束缚，打开了哑巴的口，也打开了野驴的口\BibleRef{创 16:12}，那时哈加尔的种族发出了赞美的呼喊\BibleRef{宗 2:11}。愿对祢的赞美受赞美！\par

28. 在第二十四年，愿宝库赞美圣子。宝藏对宝藏的主感到惊奇，当祂在穷人之家逐渐长大时\BibleRef{格后 8:9}，祂使自己贫穷为要使众人富足。愿祢的统治受赞美！\par

29. 在第二十五年，愿依撒格赞美圣子，因祂的良善，他在山上被从刀下救出，代替他的是一只牺牲，即那被宰杀的羔羊的预像\BibleRef{希 11:19}。那必死的人逃脱了，而那使万有复活的却死了。\EditorialNote{依撒意亚第五十三章}愿祂的祭献受赞美！\par

30. 在第二十六年，愿梅瑟与我们一同赞美祂，因为他曾害怕而从杀害他的人面前逃跑。愿他赞美那承受了长枪，在手上、脚上接受了钉子的主\BibleRef{依 49:24}。祂进入地狱并洗劫了它，然后出来。愿祢的复活受赞美！\par

31. 在第二十七年，愿善辩者赞美圣子，因为他们找不到外衣来拯救我们的案件。祂在审判厅中沉默，祂承担了我们的案件。愿光荣归于祂！\par

32. 在这一年，愿所有审判者赞美祂，他们作为义人，杀了不虔敬的人；愿他们赞美为了恶人而死的圣子，因为祂是善的。虽为义者之子，祂却丰盛地赐予他们一切美善。愿祂慈悲的肠脏受赞美！\par

33. 在第二十八年，愿所有英勇有力的人赞美圣子，因为他们并未从那掳掠我们者手中拯救我们。唯独祂堪受赞美，祂被杀害却给我们展示了生命。\BibleQuote{\BibleRef{默 5:9}} 愿祂的解救受赞美！\par

34. 在第二十九年，愿约伯与我们一同赞美祂，他为自己忍受了苦难，而我们的主为我们忍受了唾污、长矛、荆棘冠冕、鞭打、轻蔑和侮辱，是的，还有嘲弄。愿祂的慈爱受赞美！\par

35. 在第十三年，愿死者与我们一同赞美祂，因为他们得以复活；愿生者也赞美祂，因为他们已转向悔改，\BibleQuote{\BibleRef{拉 4:6}} 因为高低诸界都由祂和好为一。愿祂和祂的圣父受赞美！\par

% structure: p0278 source=h2 target=section reason=relative-heading-rank; base=h2

\section{第十四首颂歌}

（应和：——\Emph{愿祂受赞美，祂变得无限卑微，为使我们变得无限伟大}）\par

1. 论首生者的诞生，让我们在祂的庆节上讲述。——祂在自己的日子里赐下隐秘的安慰。——若不洁的\Emph{君王}在自己的宴席上，为纪念自己的日子，赐下了愤怒的礼物——即盘子中的首级，——何况蒙福者岂不更要赐福与那些在祂的庆节上歌唱赞美祂的人！\par

2. 不要将我们的守夜等同于每日的守夜。——祂的庆节，其赏报超过百倍。——因为这庆节藉守夜向睡眠宣战；它藉声音向沉默宣战；——它披戴一切福分，是庆节之首，——也是一切喜乐之首。\par

3. 今日天使和总领天使——降下，——在地上唱一首新歌。——在这奥迹中，他们降下，与守夜者一同喜乐。——当他们赞美时，亵渎却曾泛滥。——愿那诞生受赞美，看啊！世界因之——回响着赞美的颂歌。\par

4. 因为这一夜使高天的守望者与地上的守夜者联合。——守望者来临，要在受造界中设立守夜者。——看啊！守夜者成了守望者的同伴：——歌颂者成了色辣芬的伴侣。——愿那成为祢赞美的竖琴者受赞美！——愿祢的恩宠成为他的赏报。\par

5. 那么，我要歌唱并讲述首生者的诞生，——天主性如何在母腹中为自己编织了一件衣袍。——祂穿上它，在诞生时出来，在死亡时又脱下它——祂脱下过一次，又穿上过两次。——祂曾在左边穿上它，然后从那里脱下——并放在右边。\par

6. 祂住在狭窄的怀抱中，那统御万有的全能者。——当祂住在那里时，祂掌握着整个宇宙的缰绳：——祂向圣父奉献，为满全祂的旨意：——天被祂充满，以及一切受造物。——太阳进入母腹，而在高天与深渊——祂的光辉仍存留。\par

7. 祂住在一切受造物的宽阔怀抱中，——它们都过于狭窄，容不下首生者的伟大。——那么，玛利亚的怀抱如何容纳了祂？——若能容纳，是奇事；若不能容纳，则是困惑。——在一切容纳祂的怀抱中，只有一个怀抱足以容纳祂——即那生祂的至高者自己的怀抱。\par

8. 那怀抱祂的怀抱，若完全容纳了祂，——就等于那生祂的至高者的奇妙怀抱。——但谁敢说那狭窄、软弱、卑微的怀抱——等同于祂，那至高者的怀抱？——祂出于慈爱住在那里，尽管祂的本性如此伟大：——它是无界限的。\par

9. 和好的和平，被派遣到万民中！——喜乐的光明，来到忧伤者中！——沉默的有力酵母，战胜一切！——忍耐的那一位，祢已将一人又一人网入祢的网中！——有福的人，他在心中接纳了祢的喜乐——并在祢内忘记了他的叹息！\par

10. 守望者向守夜者宣告了和平。——在守夜者中，守望者传报了喜讯。——谁会在那唤醒一切受造物的夜晚睡觉呢？——因为他们传报了和平的喜讯，而那里曾有争战。——有福的人，他因沉默取悦了\Emph{神圣}威严，——而说话却激起祂的愤怒！\par

11. 守望者与守望者相混，他们因世界复活而喜乐。——恶者蒙羞，他曾为王，编织了谎言的冠冕；——并在世界中设立他的宝座，如同天主。——躺在马槽中的婴孩，将他从统治中赶出。——太阳献上朝拜，通过他的贤士们向他致敬——在他的朝拜者中，他们朝拜了祂。\par

12. 天主看见人类朝拜受造之物：——祂穿上受造的身体，为按我们的习惯俘获我们。——看啊！以这\Emph{我们}的形态，那塑造我们的医治了我们——并且在这受造的形体中，我们的创造者赐予我们生命。——祂不是用强权吸引我们：愿祂受赞美，祂以我们的样式而来——并将我们与祂联合！\par

13. 谁能不惊叹玛利亚，达味之女——怀抱婴孩，而她的童贞得以保存！——她将祂放在怀中，以歌声哄祂入睡，祂便喜乐。——天使高唱颂歌，色辣芬呼喊\Quote{圣}，——贤士们献上悦纳的礼物，——给那诞生的圣子。\par

14. 哦，无限伟大，被无限贬抑——超乎赞美，被降卑至屈辱中！——祢的慈爱将祢俯就这一切——愿祢的恩宠俯就我，虽我邪恶，也能献上赞美！——有福的人，他成为众声的源泉——全都以万有赞美祢！\par

15. 祂在地上是仆人；祂在天上是主。——高天和深渊的继承者，祂成了外方人：——人用诡计审判祂，祂却按真理审判：——他们向祂的面吐唾沫，祂却将自己的圣神吹在他们身上：——那拿着脆弱芦苇的，成了世界的支柱——世界日益衰老，便倚靠祂。\par

16. 那起来服侍祂仆人的，如今坐着接受朝拜。——经师们藐视的祂，色辣芬在祂面前呼喊\Quote{圣}。——这赞美是亚当想要偷取的。——那使他跌倒的蛇，看见他被提升到何等高度：——他粉碎了它，因它欺骗了他；厄娃的脚践踏了它——它曾将毒液注入她的耳中。\par

17. 妻子曾不生育，保留了她的果实——但玛利亚的怀抱圣洁地受孕。——她无需惊叹田地、赞赏植物——她接受了未曾借取的，并归还了它。——本性承认了自己的失败；母腹觉察到了——并归还了\Emph{本性}所不曾给予的。\par

18. 玛利亚在判断中被依撒伯尔击败。——那不能生育的恳求：那得胜的旨意，关闭了敞开的门，却打开了关闭的。——祂使已婚的子宫无子女，使童贞的子宫多结果实。——因为子民是被诅咒的无信者，祂使那已婚的在她面前被阻而不能生育。\par

19. 祂能给枯干无生命的乳房以乳汁，却使她们在青年时期干涸，在老年时期流淌——强迫并改变了本性，在合时与不合时。——万有的主改变了童贞女的本性。——因为子民是不结果实的，祂使那老妇成为少女的口。\par

20. 正如祂在出生时开始，祂继续并在死亡中完成。——祂的诞生接受了朝拜；祂的死亡偿还了债务。——当祂来到诞生时，贤士们朝拜了祂；祂又来到受难时，盗贼投靠了祂。——在祂的诞生与死亡之间，祂将世界置于中间：在出生和死亡中，祂赐予它生命。\par

21. 千千万万站立，万万奔忙。——千万与万万，不能探究那独一者：因为他们全都站住，静默事奉。——祂的宝座没有后嗣，除了那出自祂的圣子。——在沉默之中是对祂的探寻，当守卫者来探究祂时——他们达到沉默，便止住了。\par

22. 首生者进入子宫，纯洁的童贞女未受损伤。——祂在\Emph{她的}阵痛中骚动并出来，美丽的母亲因祂而苦恼。——进入时荣耀而不可见，出来时卑微而显明——因为祂进入时是天主，出来时是人。——听闻是奇事与困惑：火进入了子宫；穿上了身体而出来！\par

23. 总领天使加俾额尔称祂为\Quote{我的主}：——他称祂为\Quote{我的主}，是为了教导祂是他的主，而非同伴。——加俾额尔有弥额尔与他为伴：——圣子是众仆人的主；祂的本性\Emph{是}崇高的，如祂的名号一样。——\Emph{没有}仆人能探究祂；因为仆人越大——祂在祂的仆人之上越伟大。\par

24. 当他们站在祢面前，守卫者以赞美诗歌——他们不知道在何处能辨认祢。——他们曾在高处寻求祢，在深渊看见祢：——他们在天空中间寻找祢，在深渊中间看见祢：——他们在受朝拜者旁边辨认祢，在受造物中间找到祢：——他们下到祢那里，向祢歌唱光荣。\par

25. 祢全然奇妙，在我们寻求祢的一切方面。——祢既近又远，谁能达到祢？——没有寻求能成功，以使它的伸展能触及祢。——当它伸展以触及祢时，就被阻止而停止——它达不到祢的山；信德达到那里——还有爱德与祈祷。\par

26. 贤士们也寻求祂，在马槽找到祂时——他们以沉默而非审察来朝拜祂；——他们献上供物以代替空洞的挣扎。——你们也寻求首生者，若你们在至高之处找到祂——不要以烦扰的质问，要在祂面前打开你们的宝藏——并献上你们的工作。\par

% structure: p0306 source=h2 target=section reason=relative-heading-rank; base=h2

\section{第十五首}

应答。——\Quote{\Emph{祂在诞生中超越一切，应受赞美！}}（两次）。\par

1. 万民啊，庆祝这庆节，万节之首果；——细数曾经的苦难、创伤与痛苦——好使我们知道祂医好了什么灾祸，那被派遣的圣子。\par

应答。——\Quote{祂足以治愈我们的痛苦，应受赞美！}\par

2. 蒙救的万民啊，庆祝祂，祂在诞生中拯救一切。——连我这软弱的口舌，也因\Emph{祂的}仁慈成了竖琴。——首生者的卓越，让我们在祂的庆节中歌唱。\par

应答。——\Quote{祂使我们适于祂的庆节，应受赞美！}\par

3. 那么，有谁能赞美一位医师——除非他听闻并知道祂所治愈的痛苦是什么？——当我们的灾祸被宣告时，我们的医治者就被显扬了。\par

应答。——\Quote{祂在我们的痛苦中被高举，应受赞美！}\par

4. 受造物曾被朝拜：因为朝拜者是愚昧的——他惯于朝拜万有，却未朝拜那独一者。——因此祂仁慈地降来，打碎了奴役万有的轭。\par

应答。——\Quote{祂解除了我们的痛苦，应受赞美！}\par

5. 至高者的仁慈被启示；祂降来释放祂的受造物。——在这蒙福的月份，其中\Emph{奴隶}得释放——主经历了束缚，以召叫被束缚者得自由。\par

应答。——\Quote{祂带来了自由，应受赞美！}\par

6. 月份的主为自己选择了两个月来行事。——祂的受孕在尼散月，祂的诞生在科农月。——在尼散月祂圣化了那些受孕的；在科农月祂释放了那些出生的。\par

应答。——\Quote{祂使祂的月份喜乐，应受赞美！}\par

7. 太阳在沉默中揭示了他的朝拜者归向他的主：——身为仆人却受朝拜而非他的主，这令他痛苦。——看啊！受造物喜乐，因为造物主受朝拜。\par

应答。——\Quote{那受朝拜的婴孩，应受赞美。}\par

8. 月份戴着三个\Emph{冠冕}，在祂的凯旋中为祂加冕。——太阳\Emph{是}为祂的诞生而蒙福，为祂的复活而受渴慕——为祂的升天而蒙福；月份为祂带来了冠冕。\par

应答。——\Quote{祂在祂的月份中凯旋，应受赞美！}\par

9. 揭下面纱，喜悦你的面容吧，受造物，在我们的庆节中。——让教会以声音歌唱；天地静默！——歌唱赞美那婴孩，祂为一切带来了释放！\par

应答。——\Quote{祂废除了束缚，应受赞美！}\par

10. 当愚昧人向太阳致敬时，他们因尊敬他而羞辱了他。——但现在当众人都知道他是仆人，在他的轨道中他的主受朝拜；——众仆人都喜乐，因为他们被视为仆人。\par

应答。——\Quote{祂安排了他们的本性，应受赞美！}\par

11. 我们行了悖逆的事，成了仆人的仆人。——看啊！我们的自由强迫他，那仆人，成为我们的主：——太阳，为众人服务的仆人，我们竟使他成为众人的主。\par

应答。——\Quote{祂使我们转向祂自己，应受赞美！}\par

12. 那曾被朝拜的月亮，也因祂的诞生得了自由。——因为奇怪的是，藉着她的光，那照亮眼睛的——藉着它眼睛却变暗了，以致他们注视她如神。\par

应答。——\Quote{那光照了我们的光束，应受赞美！}\par

13. 火赞扬祢的诞生，因它引开了对自己的朝拜。——贤士们曾朝拜它：他们曾在祢面前朝拜。——他们舍弃了它而朝拜它的主；他们以火交换了那火。\par

应答。——\Quote{祂以祂的光沐浴了我们，应受赞美！}\par

14. 代替那吞噬自身躯体、自身无知的火——贤士朝拜了那赐予祂的身体作为食物的火。——活炭靠近，圣化了那不洁的双唇。\par

R.，愿祂受赞美，祂在我们内调和了祂的火！\par

15. 迷惑蒙蔽了世人，叫他们敬拜受造物：——同受造者被敬拜，万有的天主受了委屈。——应受敬拜的那位降来，到祂的诞生，将敬拜聚集于自身。\par

R.，愿祂受赞美，祂由万有所敬拜！\par

16. 全知者看见人敬拜受造之物：——祂穿上受造的身体，为照着我们的习惯俘虏我们，——并借着受造的身体，将我们引领到造物主前。\par

R.，愿祂受赞美，祂以巧计吸引了我们！\par

17. 恶者知道如何伤害我们：用光明使我们眼瞎——用财物伤害我们，用金子使我们贫穷——用雕刻匠的偶像，使我们心如铁石。\par

R.，愿祂受赞美，祂来使它软化！\par

18. 他们雕刻并竖起石头，使人绊倒其上。——他们不将石头放在大路上，叫盲人在上面绊倒：——他们称它们为神，使人睁着眼睛在其上绊倒。\par

R.，愿祂受赞美，祂揭露了他们所畏惧的偶像！\par

19. 罪恶展开翅膀，遮蔽万物——以致无人能凭自己或从上头辨明真理。——真理降入母胎，出来并滚开了谬误。\par

R.，愿祂受赞美，祂藉祂的诞生驱散了罪恶！\par

20. 因为慈悲不忍见道路受阻。——当祂为成孕而降临时，祂开了道路并使它平坦：——当祂以诞生出来时，祂踏过它并标出了里程。\par

R.，愿祢的道路平安！\par

21. 祂拣选了先知；他们为人民清除了道路：——祂派遣了宗徒；他们为万民铺平了路径。——恶者的网罗蒙羞，当脆弱的人将它们清除时。\par

R.，愿祂受赞美，祂使我们的道路平坦！\par

22. 偶像使人眼瞎，它们的雕刻师在暗中：——他们在石头上刻了眼睛，却使灵魂的眼睛变暗。——赞美祢的诞生，它开启了被弄瞎的视力。\par

R.，愿祂受赞美，祂恢复了视力！\par

23. 愿妇女赞美她，纯洁的玛利亚，——因为如同在她们的母亲厄娃身上——她们的羞辱极大——看！在她们的姊妹玛利亚身上——她们的尊荣被大大高举。\par

R.，愿祂受赞美，祂由妇女而生！\par

24. 愿万民赞美祢的诞生，因为他们获得了眼睛去看——他们的酒如何使他们眩晕；他们看见了自己的羞辱吗？——他们开始认识自己，并敬拜那拯救了他们的那一位。\par

R.，愿祂受赞美，祂教导了悔改！\par

25. 它的敬拜——人类——已散布各处：——那应受敬拜者，它却不寻求，好使敬拜归于祂。——但祂不容忍——错误的敬拜者。\par

R.，愿祂受赞美，祂降临并受敬拜！\par

26. 偶像的金子敬拜了祢，因为祢视它为施舍；它单独在生活的用途中毫无用处。——它急速奔向祢的钱囊，正如它曾急速奔向马槽。\par

R.，愿祂受赞美，受造界所爱的那一位！\par

27. 乳香敬拜了祢的诞生，它曾事奉邪魔。——它在蒸气中哀恸：它看见自己的主时便欢跃。——它不再作迷惑的香，而是在天主前的祭品！\par

R.，愿祢的诞生受敬拜！\par

28. 没药敬拜了祢，为了它自己及它的同类香膏。——那携带其香膏的手，曾膏抹可憎的偶像。——对于祢，香气是甘饴的，来自玛利亚用以膏抹祢的傅油。\par

R.，愿祢的馨香为我们是甘饴的！\par

29. 曾受敬拜的金子敬拜了祢，当贤士献上它时。——那曾在铸像中受敬拜的，向祢行了敬拜。——它同它的敬拜者一起敬拜了祢，它承认祢是应受敬拜的那一位。\par

R.，愿祂受赞美，祂为自己要求了敬拜！\par

30. 恶者和他的军队逃跑了，他常在世上欢跃。——他们在高处向他献小母牛，在园中为他宰公牛。——他吞噬了所有受造物，他的肚腹充满了猎物。\par

R.，愿祂受赞美，祂来使他呕吐！\par

31. 关于他，主曾说他从天上坠落。——那可憎者曾高举自己；从高升中他坠落了。玛利亚的脚踏碎了他，他曾用脚跟伤厄娃。\par

R.，愿祂受赞美，祂藉诞生使他屈服！\par

32. 迦勒底人遍处游行，引人入迷途：——迷惑的宣讲者在全世界蒙羞——他们蒙羞并被征服——被真理的宣讲者。\par

R.，愿那婴孩受赞美，他们宣讲了祂！\par

33. 罪恶张开了她的网，为捕鱼。——赞美归于祢的诞生，它捕获了迷惑的网。——灵魂高飞，它曾被捕获于深渊。\par

R.，愿祂受赞美，祂为我们预备了翅膀！\par

34. 祂的旨意有能力，甚至以强力拯救我们。——但既然不是强力使我们有罪，祂也不是以强力洁净我们。——恶者以诱惑奴役了我们：祢的诞生以诱惑赐给我们生命。\par

R.，愿祂受赞美，祂计划并赐给我们生命！\par

35. 受造物抱怨它们被敬拜；在沉默中它们寻求释放。——那全面解放者听见了，因祂不能容忍，祂降临——在母胎中取了奴仆的形体，出来，释放了受造界。\par

R.，愿祂受赞美，祂使受造界成为祂的收益！\par

36. 慈悲在高天被点燃，因受造界呼喊的声音：——加俾额尔被派遣；他来了，传报了祢受孕的喜讯。——当祢到了诞生时，守望者传报了祢出来的喜讯。\par

R.，愿祢因超乎万有的敬拜而受赞美！\par

37. 因为诞生的喜乐比受孕更大。——是的，是一位天使给我们带来了祢受孕的喜讯：——但在祢诞生的喜乐中，一大群守望者带来了喜讯。\par

R.，愿祢在祢的日子里的喜讯受赞美！\par

38. 在祢的日子里，我愿向祢献上光荣，哦，受敬拜的那一位！——请收下属于我的果实；并赐给我属于祢的慈悲！——因为若我内的邪恶能献礼，何况祢是善者，岂不更能赐予！\par

R.，愿祢的财富在祢的仆人中受赞美！\par

39. 祢所寻求的两件事，在祢的诞生中为我们成就了。——祢穿上了我们可见的身体；我们穿上了祢不可见的权能：——我们的身体成了祢的衣服；祢的圣神成了我们的外袍。\par

R.，愿祂受赞美，祂被装饰并装饰了我们！\par

40. 高天与深渊惊异，因祢的降生制伏了叛逆者。——因为我们给了祢人质，祢赐给了我们\OriginalTerm{护慰者}{Paraclete}。——人质从我们这里上升时，万军之统帅降临到我们当中。\par

R.，愿那取去又降下的受赞颂！\par

41. 来吧，众口，倾流吧，如同诸水，成为声音的源泉！愿圣神降临——藉着我们众人向圣父歌唱光荣，祂藉着自己的圣子救赎了我们！\par

\Emph{R.，愿祂在祂的降生中超越万有受赞颂！}\par

% structure: p0390 source=h2 target=section reason=relative-heading-rank; base=h2

\section{颂歌 16}

（应答。——\Emph{愿我们众人向祢全体献上光荣！}（两次））\par

1. 有哪个属血肉的人，能述说那赐予万有生命者——祂舍弃了尊威的高天，屈尊就卑？——\Emph{祢}在祢的降生中高举万有，求祢高举我软弱的心智——以述说祢的降生；不是要让我探究祢的尊威，——而是要让我传扬祢的恩宠。\par

R.，愿那在祂的讲论中隐藏又启示的受赞颂！\par

2. 这是一个伟大的奇迹：圣子完全居住在一个身体内——完全住在其中，这身体便足以容纳祂；祂住在其中，却不受其限制。——祂的旨意完全在其中；祂的界限完全达到圣父。——谁能足能说明，祂虽完全住在一个身体内——同样也完全住在万有中？\par

R.，愿那虽无界限却被限定的受赞颂！\par

3. 祢的尊威向我们隐藏；祢的恩宠向我们彰显。——主，关于祢的尊威，我将静默；关于祢的恩宠，我必讲述。——祢的恩宠依附于祢，使祢俯就我们的卑微：——祢的恩宠使祢成为婴孩；祢的恩宠使祢成为人：——它限制了祢，又扩展了祢的尊威。\par

R.，愿那变小又变大的大能受赞颂！\par

4. 光荣归于那成为卑微者，祂本性原是崇高！——祂因着爱成为\Person{玛利亚}的长子，本人却是神性的首生者。——祂在名分上成为\Person{若瑟}的后裔，本人却是至高者的后裔。——祂凭自己的旨意成为人，按本性却是天主。——愿祢的旨意和祢的本性受赞颂！\par

R.，愿那披上我们形象的祢的荣耀受赞颂！\par

5. 是的，主，祢的降生成了\Emph{一切}受造物的母亲；因为它重新经历产痛，生下了那生养祢的人类。祢由它肉体而生；它由祢属灵而生。——祢来到世间降生，\Emph{是为了}人类能按祢的形象重生。——祢的降生成了万有重生的创始者。\par

R.，愿那成为青年并赐予万有青春者受赞颂！\par

6. 当人类的希望破灭时，祢的降生增加了希望。——天上诸神向人类传报了希望的喜讯。——那斩断我们希望的\Person{撒殚}，亲手斩断了它自己的希望。——当他看见希望增加时：祢的降生成了绝望者——一个充满希望的泉源。\par

R.，愿那传报希望者受赞颂！\par

7. 祢降生的日子如同祢，因它如祢一般可喜爱、可渴望。——我们这些未见过祢降生的人，以及它如同在其时代中的火焰，——在这个祢的日子，我们看见祢，正如祢曾是婴孩——为众人所爱，看！在祢内众教会喜乐——祢的日子装饰并受装饰。\par

R.，愿那为我们预定的祢的日子受赞颂！\par

8. 祢的日子赐给了我们一份恩赐，圣父没有其他恩赐像它。——祂没有派遣\OriginalTerm{色辣芬}{Seraphim}给我们，\OriginalTerm{革鲁宾}{Cherubim}也没有降临到我们中间——没有\OriginalTerm{守者}{Watchers}或\OriginalTerm{服役者}{Ministers}降临，只有他们侍奉的\OriginalTerm{首生者}{Firstborn}来临。——谁能足能感恩，那不可测量的尊威——竟躺在卑贱的马槽中！\par

R.，愿那赐给我们祂所赢得之物的受赞颂！\par

9. 祢的降生使那一代喜乐，我们的世代因祢的日子喜乐：那一代有双重的福乐，因他们看见了祢的降生和祢的日子：——后世者的福乐较少，因他们只见祢降生的日子。——但因那时的人怀疑，后世者的福乐更大——他们虽未看见祢，却相信了祢。\par

R.，愿那加给我们的祢的福乐受赞颂！\par

10. 贤士们从远方尊崇；经师们在近旁低语——先知传达了他的讯息，\Person{黑落德}传达了他的愤怒——经师们显示了他们的学问，贤士们显示了他们的供物。这是一个奇迹：对祂——那婴孩，祂家里的人携刀匆匆前来，而外邦人带着供物前来。\par

R.，愿那激动万有的祢的降生受赞颂！\par

11. \Person{玛利亚}的怀抱令我惊奇，它竟足以容纳祢，主，并拥抱祢。——所有受造物都太小，不能隐藏祢的尊威；——天地太窄，不能如同翅膀覆盖祢的神性。——大地的怀抱对祢太小；\Person{玛利亚}的怀抱对祢足够大。——祂住在怀抱中，并在她的怀抱中医治。\par

R.，\par

12. 祂被简陋地包裹在襁褓里，供物被献予祂。——祂在青年时穿上衣服，从它们中发出帮助：祂穿上圣洗的水，从它们中射出光芒：——祂在死亡中穿上细麻布，在它们中显出了胜利；伴随着祂的卑微，祂的升高。\par

R.，愿那将祂的光荣与祂的苦难结合者受赞颂！\par

13. 这些都是衣袍的改变，是慈悲脱下又穿上的——当祂努力给\Person{亚当}穿上那祂已脱下的荣耀时。——祂被裹在襁褓中，如同\Person{亚当}被树叶包裹；又穿上衣服代替皮衣。——祂为\Person{亚当}的罪受洗，为\Person{亚当}的死被埋葬：——祂复活并高举\Person{亚当}进入荣耀。\par

R.，愿那降临、给他穿上又上升者受赞颂！\par

14. 虽然祢的降生，对\Person{亚当}的子孙如同对\Person{亚当}已经足够；——大能者啊，祢成了婴孩，在祢的降生中祢重新生了我！——纯洁者啊，祢受了洗，愿祢的圣洗洗去我们的污秽——永生者啊，祢被埋葬，愿我们在祢的死亡中获得生命！——我要在充满万有的祂内赞美你们全体。\par

\Emph{R.，愿我们众人向祢全体献上光荣！}\par

% structure: p0420 source=h2 target=section reason=relative-heading-rank; base=h2

\section{颂歌 17}

（应答。——\Emph{在祢降生的这一天，愿赞美从众口归于祢！}）\par

1. 因祢的降生，婴儿被杀害，祢是万有的生命赐予者——但因那被杀者是君王，我们主、诸国的主，——暴君诡诈地，为了祂献上被杀的人质，——穿着祂被杀之奥秘的衣袍：天上的万军领受了——地上所献的人质。\par

R.，愿那尊崇祂的君王受赞颂！\par

2. 达味家族的所有君王，一代传一代，拖拉传递——达味之子的宝座和王冠，如同保管存款。——在一人身上，他们达到了界限和极限，当祂，万有之主，来到时——便从他们手中夺走了一切，切断了万有的传递……\par

应：愿那穿戴属于自己者受赞美！\par

3. 鸽子在白冷哀鸣，因为蛇毁灭了它们的后裔。——鹰飞往埃及，要下去领受恩许。——埃及因祂的到来而欢喜，因祂带来了丰盈，以偿还债务——这债务是若瑟的子孙所未能清偿的。在若瑟的子孙中，祂劳苦并偿还了——若瑟子孙的债务。\par

应：愿那从埃及召叫祂者受赞美！\par

4. 经师们每日诵读，说星要从雅各伯升起。——因为选民既是\Emph{是}声音又是诵读，对外邦人则是\Emph{星的}升起与解释：——书卷为他们，事实为我们；枝杈为他们，果实为我们。——经师在书写之物中诵读；贤士在成就之事中看见那被诵读之物的闪耀。\par

应：愿那为我们增添他们书卷者受赞美！\par

5. 谁能述说那退隐与显现——那在献礼者前行、发光的星？——它显现并宣告了冠冕；它隐藏并遮盖了\Emph{祂的}身体。——它为了圣子有双重作用：先驱与护卫；——它护卫了祂的身体，它宣告了祂的冠冕。\par

应：愿那赐智慧给宣扬祂者受赞美！\par

6. 暴君注视着贤士，当他们问：\Quote{君王之子在哪里？}——他心中忧郁，却强装欢颜。——他随羊群派出豺狼，要杀害天主的羔羊。——羔羊下到埃及，好从那里审判他们——从他曾拯救他们的地方。\par

应：愿那再次制服他们的那一位受赞美！\par

7. 贤士向暴君宣告：\Quote{当你的仆从加入我们时——明亮的星退隐了，是的，道路也隐藏了。}——有福者不知道，君王已派遣了凶恶的\Emph{仇敌}——\Emph{假装}朝拜者的杀手，去毁灭那甜美的果实——吃了这苦果的，反变为甘甜。\par

应：光荣归于祢，生命的良药！\par

8. 在那里，贤士接受了命令，要去寻找祂。——经上记载他们看见了那明亮的星，便大喜乐。——\Emph{由此}可知它曾退隐；所以他们因见到它而喜乐。——它隐藏起来，阻止了凶手；它升起，召唤了朝拜者；——它倾覆了一部分，召唤了一部分。\par

应：愿那在双方都获胜者受赞美！\par

9. 那可憎的屠杀婴孩者，他如何忽视了那婴孩？——正义阻止了他，使他以为贤士会回到他那里。——当他留下等候捕捉时，被朝拜者与祂的朝拜者——一切都从他手中逃脱，供物与朝拜者都飞走了——从暴君那里去了君王之子那里。\par

应：光荣归于洞悉一切计谋的那一位！\par

10. 无瑕的贤士在睡眠中，在床上默想：——睡眠成了明镜，梦境如光升于其上。——他们看见了凶手，颤抖起来，因他的诡计与刀剑闪烁。——他教人诡诈，他磨快了刀刃：——警醒者教导了沉睡者。\par

应：愿那赐明智给朴实者的那一位受赞美！\par

11. 相信的朴实者已经知道基督的两次来临：——但愚昧的经师甚至连一次来临也未察觉。——然而外邦人在第一次来临中有了生命，还要在第二次来临中复活。——那心思被蒙蔽的选民，第一次来临驱散了他们——第二次来临将涂抹他们的记忆。\par

应：愿那已来并要求的那位君王受赞美！\par

12. 当救主如盲者般升起时，太阳显露出光芒——而他们却披上黑暗：光明放射出光华——祂带来了众星之子，以显明黑暗之子。——看啊！星在你们中间，但你们的眼上有帕子。\par

应：光荣归于祢，新生的太阳！\par

13. 先知们宣告了祂的诞生，但未明示其时间。——祂派遣了贤士，他们前来并显示了时间。——但那些揭示时间的贤士，并未明示那孩子\Emph{是谁}。——一颗灿烂光华的星，在运行中显示了那孩子是谁——\Emph{祂的}世系是\Emph{何等}辉煌。\par

应：愿那藉着他们全被指出的那一位受赞美！\par

14. 他们轻视依撒意亚的号角，它吹响了祂无玷的受孕；——他们止息了圣咏的琴弦，它歌唱了祂的司祭职；——他们缄默了圣神的竖琴，它又歌唱了祂的王国；——在深沉的寂静中，他们封闭了那伟大的诞生，它连结起上界与下界的欢呼。\par

应：愿那在寂静中显现的那一位受赞美！\par

15. 祂的声音是秘密的钥匙，开启了贤士的口。——当宣讲者在犹大沉默时，他们使自己的声音响彻受造界；——那些远方来的人接受并带走了那些人所轻视的福音。——轻蔑者开始从陌生人那里听见他们\Emph{自己}的命令，这些陌生人呼喊出\Emph{达味之子的名号}。\par

应：愿那藉着我们的声音使他们沉默的那一位受赞美！\par

16. 当选民轻视供物，不将供物献给祂——君王之子时，祂派遣了先驱到外邦人中，使他们带着供物前来：——但祂并未使他们全部前来，因为白冷狭窄的怀抱不能容下他们；但圣教会的怀抱——扩大了自己，容纳了她的子女。\par

应：愿那使不生育者结果实的那一位受赞美！\par

17. 白冷的刽子手割倒了嫩花，好让嫩苗在其中灭亡，嫩苗中隐藏着生命之粮。——但那有生命的麦穗逃脱了，好来到收割时的禾捆中：——那幼时逃脱的葡萄，把自己交给了压榨——好使它的酒赐予灵魂生命。\par

应：光荣归于祢，生命的宝库！\par

18. 凶手们进入了一个乐园，充满了嫩果：——他们从枝头摇落花朵，摧毁了蓓蕾与嫩芽——他无意间献上了无瑕的祭品，那迫害者。——祸哉他，但福哉他们！白冷首先将童贞的果实献给了圣者。\par

应：愿那接受初果的那一位受赞美！\par

19. 经师们因嫉妒而沉默，法利塞人因妒忌而静默。——石头般的人呼喊赞美，他们本有石心。——他们在那被弃绝却成了角石的石头前鼓掌。——石头因那石头成了血肉，得了口舌说话；石头借着那石头呼喊。\par

应答：愿您的诞生受赞美，它使石头呼喊！\par

20. 经上记载的星，远处的外邦人看见了——为使近处的百姓蒙羞；受教而自高自大的百姓！竟被外邦人反过来教导，他们如何、在何处看见了那异象——巴郎所宣讲的；一个外人传扬它——外人看见了它。\par

应答：愿那位激发祂自己家人嫉妒者受赞美！\par

21. 愿我的恳求接近你的门，我的贫穷接近你的宝库！——请毫无尺度地赐予我，主，如同天主赐予人！——虽然你作为受赞颂者之子倍增恩赐，虽你作为君王之子增添它们；——我虽忘恩如\Emph{一切}尘土造物，亚当之子\Emph{也是如此}，——受赞颂者之子\Emph{也是如此}。\par

\Emph{应答：愿您如同您父一样得赞美！}\par

% structure: p0464 source=h2 target=section reason=relative-heading-rank; base=h2

\section{赞美诗 18}

应答——\Emph{愿派祂来的那位受赞美！}（重复）\par

1. 你是有福的，教会！看，在你内有——伟大庆节君王庆节的声音！——熙雍荒凉，她的城门干渴——而被节庆遗弃。——你的城门敞开却不饱满，你的厅堂扩大却仍不足！——在你中间，看哪！有列国呼喊并使百姓静默的声音。\par

2. 你是有福的，教会！在你的节庆中，——守望者在你的欢庆中喜乐！——因为一夜之间，守望者在地上——曾赞美，这地上中止并拒绝赞美。——你的声音有福，它们被播种、收割——并积存在天上的仓库！——你的口是香炉，你的声音如馨香，在你的节庆中散发香气。\par

3. 你是有福的，教会！因为在这庆节中，——所有祭品都带给你。——贤士们曾在叛逆者中，将祭品献给真理。——你的居所有福，祂屈尊并居住其中，君王之子受礼物敬拜！——西边来的黄金，东边来的香料——在你的节庆中被献上。\par

4. 你是有福的，教会！因为你那里没有——暴君君王杀害婴孩！因为他在白冷滥杀小孩子——为要杀死那赐生命给众人的孩子。——你的孩童有福，他们被君王嫉妒并敬拜，——因为那些应许为你的敬拜，——东方的冠冕：那践踏你亲爱者的人，将被你的爱子践踏。\par

5. 你是有福的，教会！看，在你身上——依撒意亚也在他的预言中欢欣——看，童贞女将怀孕生子——祂的名字是伟大的奥秘！——在教会中启示的解释！——两个名字结合成为一——厄玛奴耳，——天主永远与你同在，祂把你与祂的肢体结合！\par

6. 你是有福的，教会！在米该亚呼喊中——\Quote{一位牧者将出自厄弗辣大}：——因为祂来到白冷，为从那里——取来叶瑟的棍杖，统治列国。——你的羔羊有福，它们被印上祂的印记，——你的羊群受祂的剑保护！——你，教会，——是永久的白冷——因为在你内是生命之粮！\par

7. 你是有福的，教会！看，在你内——达尼尔，那蒙爱的人，也喜乐，——他预言受颂扬的默西亚将被杀，——圣城在祂被杀时成为荒芜！——那被弃绝而不悔改的百姓有祸了！——被召而不转离的列国有福！——受邀的客人拒绝——别人取代他们享用了筵席。\par

8. 你是有福的，教会！因为看，在你的琴上，——达味君王在你内咏唱圣咏！他在圣神中歌唱祂\Quote{祢是我的儿子，我——今日生了祢}在圣洁的光荣中。——你的耳朵有福，它们被洁净去听！——在祂的日子，醒着如同祂的身体，呼求祂——受熙雍教导——它使祂的庆节忧伤；使那使你喜乐者喜乐。\par

9. 你是有福的，教会！因为一切节庆——都已从熙雍逃逸，栖息于你！——疲倦的先知们在你内找到了安息——脱离他们在犹大所受的劳苦和羞辱。——你的书卷有福，在圣殿中展开，——你的节庆在圣所中举行！——熙雍被遗弃——看！今天列国在你的节庆中呼喊。\par

10. 你是有福的，教会！因着十个祝福，——我们的主将之作为完全的奥秘赐下：——因为一切数字都系于十，所以你因十个祝福得以完善。——你的冠冕有福，它们被编成——每一冠冕中混合所有祝福。——有福者！——戴上一切祝福，也请将你的祝福赐给我！\par

11. 你是有福的，厄弗辣大，君王的母亲！因为从你出了冠冕之主！——米该亚告知你祂从永远就有，祂的时间界限不可测度。——你的眼睛有福，它们最先认出祂！——祂认为你配得看见祂显现之时——祝福的元首，——喜乐的开端，你首先领受。\par

12. 你是有福的，白冷！因为村镇嫉妒你——和坚城！——正如它们嫉妒你，同样妇女们嫉妒玛利亚，——和王子们的女儿童贞女。——那位愿居留其中的少女有福，——那位愿寄居其中的城邑有福——一个贫穷的少女，——一座小城，祂拣选使自己卑微。\par

13. 你是有福的，白冷！因为在你内是为祂作了开始——祂，圣子，从永远就在圣父内！——这难以领悟：祂在时间之前就存在——却在你内使自己服从时间。——你的耳朵有福，因为在你内首先听见了——天主的羔羊在你内欢乐的呼声！——你的马槽虽窄小——祂却广布四方，受万物朝拜。\par

14. 你\Emph{是}有福的，玛利亚，因你的名字——因你的孩子而伟大崇高！——你能述说那伟大者如何、多久——又在何处居于你这狭小之处。——你有福的口，曾赞美而不追问——你的舌，曾颂扬而不质疑！——因为祂的母亲尚且对祂存疑——即使她已在\Emph{腹中}怀了祂；那么谁能完全\Emph{理解}祂呢？\par

15. 女人，你无人所识——我们如何得见你所生的儿子？——因为双目不足以站立——面对那荣耀的光辉变形，那光辉在祂身上。——因为火舌居于祂内——祂藉升天派遣火舌。——愿每一舌受警戒——我们的追问\Emph{如同}碎秸，我们的审视\Emph{如同}火焰。\par

16. 那位司铎\Emph{是}有福的，他在圣所中——向圣父献上圣父之子——从我们树上摘取的果实，虽然它全然属乎\Emph{神圣}尊威！——那被祝圣而献上祂的手是有福的！——那亲吻祂而竭尽的嘴唇！——圣殿中的圣神——渴望祂的拥抱；在祂被钉十字架时，撕裂了\Emph{帐幔}而出。\par

17. 总领天使向你问候——作为圣洁的凭质——大地对他成了新天，——当守望者降下并在其上歌唱荣耀。——至高者的众子环绕你的居所——因为君王之子居于你内。——你在下的居所——因守望者群而变得如同上界的天。\par

% structure: p0483 source=h2 target=section reason=relative-heading-rank; base=h2

\section{第十九首}

（应答句：——愿你的降生受赞颂，祂使万物喜乐！）\par

1. 第一年，我们的救主诞生之年——是祝福的源头，生命的基础——因为由它而生——多种胜利，一切援助的总和：——正如\Quote{起初}的第一日，——万有的伟大支柱——承载创造的建筑；——同样，首生子的年岁\Emph{承载}对人的援助。\par

2. 第二年，我们救主诞生的第二年——贤士欢欣，法利塞人悲哀：——宝库被开启——君王们赶来，婴孩被杀害。——因为在其中，在伯利恒献上了——宝贵而可畏的祭品——因为爱献上了黄金——仇恨却以刀剑献上了婴孩。\par

3. 那普照之日，因祂的诞生而喜乐——光辉的柱石，以其光芒驱散——黑暗的作为。按照那日之原型，光被造，——分裂了那覆盖在创造荣美之上的黑暗；——我们救主诞生的光辉——进来驱散那\Emph{在}人心上的黑暗。\par

4. 第一日，源头和开始——安排根苗，使万物生长。——我们救主的日子——远被颂扬于其上，是一棵栽种在世上的树。——因为祂的死亡如同地下的根；祂的复活如同天上的头；四周祂的言语如同枝条\Emph{展开}；同样，祂的身体作为果实给食者。\par

5. 愿第二日歌颂第二位儿子的诞生——以及祂的声音，这声音首先——命令穹苍，穹苍便立定——分开\Emph{上面的}水，聚集\Emph{下面的}海。——祂分开水与水，也分开自己与守望者，降下到人间。——因为那水祂曾命令聚集，——祂劈开生命之泉，赐人饮用。\par

6. 愿第三日用各种赞歌编织——诗篇的冠冕，以同一声音献上——为祂的诞生，祂在第三日使——嫩芽和花朵生长。——但现在，那赐予一切生长的——已降下，成为全然神圣的花朵；从干渴的大地萌发而上——为要装饰并加冕得胜者。\par

7. 愿第四日赞美那在造物中——于第四日创造两个光体的——祂的诞生，——愚人敬拜这些光体，却无视觉且盲目。——光体之主已降下——从母腹中如太阳照耀我们。——祂的光辉开启了盲人的眼：——祂的光芒照亮了迷途者。\par

8. 愿第五日颂扬祂，祂在第五日创造了——爬行之物和大龙——其中一类是蛇。——它用诡计欺骗了我们的母亲，一个无谋的少女。——那欺骗者曾戏弄少女，——却被鸽子揭穿为谎言，——鸽子从童贞的胸怀跃出，前来——智者踏碎了那狡猾者。\par

9. 愿第六日颂扬祂，祂在黄昏之日创造了亚当，撒旦嫉妒他；装作朋友——安慰他，却在食物中献上毒药。——生命的药到达他们二者——披上身体，接近双方。——必死的人尝了祂，因祂而活——那吞吃祂的吞噬者却落空。\par

10. 愿第七日祝圣那圣者——祂祝圣安息日，赐所有活物安息。——那从不疲倦的福者——眷顾人类，也眷顾野兽。——当自由落在轭下——祂来至降生，成为奴仆以使之自由：——在审判厅中被仆人掌掴——祂作为主，折断了那\Emph{在}自由人身上的轭。\par

11. 愿第八日，它曾为希伯来人行割损，——赞美那命令其同名者若苏厄——用火石为百姓行割损的祂，百姓身体受了割损，内心却是亵渎的。——看！在第八日，作为婴孩，——祂来到割损，祂却割损万有。——虽然亚巴郎的记号\Emph{在}祂肉身上——但熙雍盲目的女儿玷污了它。\par

12. 愿第十日按其顺序歌唱赞美。——因为\Emph{天主}是耶稣（美名！）的\Emph{第一}个字母，在数字中是十。——祂\Emph{是}如同羔羊，将数字回转。——因为当数目上到十，它便回转从一重新开始。哦，在耶稣内的伟大奥迹！祂的力量使整个创造回转！\par

13. 在祂洁净的日子，那全洁的首生子——洁净了首生子的洁净，并在\Emph{圣殿中}被献上：——献祭之主需要献祭，——以献上鸟类。——在祂的诞生中，预象得以成全，——在祂的洁净和割损中，寓意得以圆满。——祂降临，偿清债务——在复活中祂上升，并降下宝藏。\par

\SourceRecord{Translated by J.B. Morris (Hymn nos. 1-13) and A. Edward Johnston (Hymn nos. 14-19).；Nicene and Post-Nicene Fathers, Second Series；Vol. 13.；Edited by Philip Schaff and Henry Wace.；Buffalo, NY: Christian Literature Publishing Co.,；1898.；<http://www.newadvent.org/fathers/3703.htm>.}
